
\documentclass[12pt, a4paper, twoside]{article}

%--Packages--
\usepackage[english]{babel}
\usepackage[utf8]{inputenc}
\usepackage[T1]{fontenc}
%\usepackage{showkeys}

\usepackage{latexsym,amsmath,amsthm,amssymb,bm,lmodern,titlesec,xcolor,graphicx,setspace,pdfpages,
geometry,marginnote,caption,array,sectsty,siunitx,eurosym,colortbl,float,wrapfig,vmargin,tikz,hyperref}
\usepackage[nottoc]{tocbibind}
\usepackage{tikz-cd}
%\usepackage{}
\usepackage{lipsum}
\usepackage{url}

\setpapersize{A4}
\setmarginsrb{30mm}{20mm}{30mm}{30mm}%left top right bottom
                 {8mm}{6mm}{10mm}{10mm}%headheight headsep footheight footskip

% Some useful colors
\definecolor{Rosso}{RGB}{232,61,61}
\definecolor{Viola}{RGB}{188,30,188}
\definecolor{Celeste}{RGB}{34,186,211}
\definecolor{Arancione}{RGB}{232,181,19}
\definecolor{Blu}{RGB}{0, 44, 92}
\definecolor{Azzurro}{RGB}{57,186,238}


\hypersetup{colorlinks=true,linkcolor=Blu,allbordercolors=white, urlcolor=blue, citecolor=Viola}
\usepackage{enumitem} 
\renewcommand\labelitemi{\textbullet} 


% math environments
%  we use \newtheorem to specify our theorem-like environments
   % (theorem, definition, etc.) and how to display and number them.
   %
   % Note: The \theoremstyle declarations affect the appearance of the
   % Theorems, Definitions, etc.
   % Note 2: If you want swedish names you have to go and change the argument in brackets e.g. \newtheorem{thm}[defi]{Theorem} should become \newtheorem{thm}[defi]{Sats}
\theoremstyle{definition}
\newtheorem{definition}{Definition}[section]

\newtheorem{hottake}{Hot Take}[section]

\newtheorem*{problem}{Problem}
\newtheorem{conjecture}[definition]{Conjecture}

%\theoremstyle{plain}
\newtheorem{theorem}[definition]{Theorem}
\newtheorem{proposition}[definition]{Proposition}
\newtheorem{corollary}[definition]{Corollary}
\newtheorem{lemma}[definition]{Lemma}
%\newtheorem*{remark}{Remark}%[definition]{Remark}
%\newtheorem{example}[definition]{Example}
%\newtheorem*{notation}{Notation}
%\newtheorem*{note}{Note}

\theoremstyle{remark}
\newtheorem*{question}{Question}
\newtheorem{remark}[definition]{Remark}
\newtheorem{example}[definition]{Example}
\newtheorem*{notation}{Notation}
\newtheorem*{note}{Note}

%\onehalfspacing
\numberwithin{equation}{section}

%--Heading of  section--
\titleformat{\section}[display]
  {\normalfont\bfseries\color{Blu}}
  {\filleft
    \begin{tikzpicture}
    \node[
      outer sep=0pt,
      text width=2.5cm,
      minimum height=3cm,
      fill=white,
      font=\color{Blu}\fontsize{80}{0}\selectfont,
      align=center
      ] (num) {\thesection};
    \node[
      rotate=90,
      anchor=south,
      font=\color{black}\large\normalfont
      ] at ([xshift=-5pt]num.west) {\Huge\textsc{}};  
    \end{tikzpicture}%
  }
  {10pt}
  {\titlerule\vskip4pt\Huge\sc\sffamily}

%--Heading of subsection--
\newcommand*\numb[1]{%
\begin{tikzpicture}[baseline=-0.7ex]
\node[
  outer sep=0pt,
      text width=0.6cm,
      minimum height=0.6cm,
      fill=Blu,
      font=\color{white}\fontsize{12}{20}\selectfont,
      align=center
      ] (num) {\thesubsection};
\end{tikzpicture}%
}
\titleformat{\subsection}
  {\normalfont\color{Blu}\large\sc\sffamily}{\numb{\thesubsection}}{0.8em}{}

%--Heading subsubsection--
\newcommand*\stocaz[1]{%
\begin{tikzpicture}[baseline=-0.7ex]
\node[
  outer sep=0pt,
      text width=0.8cm,
      minimum height=0.5cm,
      fill=Blu,
      font=\color{white}\fontsize{10}{20}\selectfont,
      align=center
      ] (num) {\thesubsubsection};
\end{tikzpicture}%
}
\titleformat{\subsubsection}
  {\normalfont\color{Blu}\bfseries\sc\sffamily}{\stocaz{\thesubsubsection}}{0.8em}{}
  
% Blackboard Boldface
\newcommand{\NN}{\mathbb{N}}
\newcommand{\ZZ}{\mathbb{Z}}
\newcommand{\QQ}{\mathbb{Q}}
\newcommand{\RR}{\mathbb{R}}
\newcommand{\CC}{\mathbb{C}}
\newcommand{\KK}{\mathbb{K}}
\newcommand{\FF}{\mathbb{F}}

\newcommand{\1}{\mathbf{1}}
\newcommand{\0}{\mathbf{0}}

% Text
\newcommand{\GL}{\text{GL}}
\newcommand{\SL}{\text{SL}}
\newcommand{\SymP}{\text{Sym}}
\newcommand{\AltP}{{\bigwedge}}
\newcommand{\Hom}{\text{Hom}}
\newcommand{\End}{\text{End}}
\newcommand{\SO}{\text{SO}}
\newcommand{\SU}{\text{SU}}
\newcommand{\id}{\text{id}}
\newcommand{\Stab}{\text{Stab}}
\newcommand{\Cent}{\text{Cent}}
\newcommand{\Norm}{\text{Norm}}
\newcommand{\Type}{\text{Type}}
\newcommand{\Dih}{\text{D}}
\newcommand{\Sym}{\mathcal{S}} %{\mathfrak{S}} % \text{S}
%\newcommand{\Alt}{\text{A}}
\newcommand{\Klein}{\text{V}}
\newcommand{\Cyc}{\mathcal{C}}%{\text{C}}
\newcommand{\sgn}{\text{sgn}}
\newcommand{\Span}{\text{Span}}
\newcommand{\im}{\text{im }}
\newcommand{\Tr}{\text{Tr }}
\newcommand{\Reg}{\text{Reg}}
\newcommand{\Triv}{\text{Triv}}
\newcommand{\Alt}{\text{Alt}}
\newcommand{\Perm}{\text{Perm}}
\newcommand{\Stan}{\text{Stan}}
% Vectors
\newcommand{\bas}{\mathbf{\hat{e}}}
\newcommand{\fbas}{\mathbf{\hat{f}}}

\newcommand{\ihat}{\mathbf{\hat{\imath}}}
\newcommand{\jhat}{\mathbf{\hat{\jmath}}}
\newcommand{\ghat}{\mathbf{\hat{g}}}
\newcommand{\hhat}{\mathbf{\hat{h}}}
\newcommand{\vhat}{\mathbf{\hat{v}}}
\newcommand{\what}{\mathbf{\hat{w}}}
\newcommand{\uhat}{\mathbf{\hat{u}}}
\newcommand{\xhat}{\mathbf{\hat{x}}}

\newcommand{\vvec}{\mathbf{v}}
\newcommand{\uvec}{\mathbf{u}}
\newcommand{\wvec}{\mathbf{w}}


\newcommand{\Ss}{\scriptstyle}

%\newcommand{\barpi}{\overline{\pi}}
\newcommand{\barpi}{\bar{\pi}}

\allowdisplaybreaks
\begin{document}

% This includes the front page (you need to compile the cover.tex separately and have the .pdf available in the same folder as this .tex file).
\setlength{\voffset}{0cm}\setlength{\hoffset}{0cm}\includepdf[pages={1,2,3,4}]{cover}\setlength{\voffset}{-2.54cm}\setlength{\hoffset}{-2.54cm}



%\clearpage{\thispagestyle{empty}\cleardoublepage}

%\begin{abstract}
%    Your summary goes here
%\end{abstract}

%create the table of contents
\clearpage{\thispagestyle{empty}\cleardoublepage}
{\hypersetup{linkcolor=black}
\tableofcontents
}
\onehalfspacing
\clearpage{\thispagestyle{empty}}
\begin{center}
	\subsubsection*{Abstract}
	\addcontentsline{toc}{section}{Abstract}
\end{center}

	This text serves as an introduction to representation theory of finite groups, beginning with a background in group theory and linear algebra before formally defining representations, some representations are shown to be irreducible, while others are composed as a direct sum of the irreducibles (Maschke's Theorem~\ref{thm:maschkes}), culminating in the proof of Schur's Lemma~\ref{thm:schur}, which ensure the uniqueness of these compositions. Finally, character theory is introduced, simplifying representation theory by focusing on the trace of matrices, providing systematic methods to identify and decompose any given representation. Examples to demonstrate the these theoretical concepts are some cyclic and symmetric groups of small order. In particular, $\Cyc_3$, $\Cyc_4$, $\Cyc_5$, $\Sym_3$, and $\Sym_4$ are examined and completely decomposed and the results are presented in Tables~\ref{table:Cyc3},~\ref{tbl:cyc4},~\ref{table:Cyc5},~\ref{table:completecharS3}~and~\ref{table:completecharS4}.
		
\begin{center}
	\item\subsubsection*{Sammanfattning}
	\addcontentsline{toc}{section}{Sammanfattning}
	
\end{center}

	Den här texten ger en introduktion till  ändliga gruppers representationsteori. Utifrån en inledning med ett par bakgrundsämnen inom gruppteori och linjär algebra definieras representationer formellt. Vissa representationer visar sig vara oreducerbara, medan andra visar sig vara den direkta summan av oreducerbara representationer (Maschkes Sats~\ref{thm:maschkes}), vilket kulminerar i beviset av  Schurs Lemma~\ref{thm:schur} som visar att dessa sammansättningar är unika. Slutligen introduceras karaktärsteorin som förenklar arbetet med att analysera representationer genom att undersöka matrisers spår. Karaktärsteorin ger systematiska metoder för att identifiera och sönderdela varje given representation. För att demonstrera dessa teoretiska koncept ges några cykliska och symmetriska grupper som exempel, särskilt så undersöks och sönderdelas $\Cyc_3$, $\Cyc_4$, $\Cyc_5$, $\Sym_3$ och $\Sym_4$ helt och resultaten presenteras i Tabellerna~\ref{table:Cyc3},~\ref{tbl:cyc4},~\ref{table:Cyc5},~\ref{table:completecharS3}~och~\ref{table:completecharS4}.
	
%	Sammanfattningen på svenska hamnar här. %\lipsum[2][1-10]


\begin{center}
	\item\subsubsection*{Acknowledgements}
	\addcontentsline{toc}{section}{Acknowledgements}
\end{center}
	
%	\textit{NOT DONE}
	
	I would like to thank my supervisor Sofia Tirabassi for guiding me through the construction of this text. I also would like to thank Stefan Sigurdsson for proofreading the text.
	
%	Artificial intelligence (AI) services were not used in the process of planning, writing or typesetting this text. 
%	This text was typed with {\LaTeX} in TexStudio versions $\leq4.7.3$.
	%\lipsum*[3][1-10]



\clearpage{\thispagestyle{empty}}
\section{Introduction}

	The main purpose of this text is to study linear representations of groups, or in other words replacing abstract groups with sets of matrices behaving the same way as the elements of the group and replacing the group operation with matrix multiplication. This is the intuition behind a homomorphism.
	
	The reader is assumed to be familiar with first-year linear algebra and the definitions and basic concepts of vector spaces and groups.
	
	Only representations of finite groups in finite dimensional vector spaces will be treated. As \cite[Sect.1.1.]{Serre} notes, we are usually interested in a finite number of elements of a vector space, and then we could span a finite subspace with those elements. The restriction to finite groups is more severe however, as some results in this text does not apply on infinite groups, or even compact infinite groups (see \cite{FultonHarris, Serre}). 
	
	Calculations performed in this text are my own and the examples are also my own, except where otherwise noted.
	
	As examples, we study the abelian cyclic groups and the symmetric groups of small order. The cyclic group of order 3, $\Cyc_3$ and the symmetric group of order $\Sym_3$ are used extensively as examples.
	
	Section 2 introduces some background topics from group theory and linear algebra required before introducing representation theory. Especially the symmetric group, recollections of linear algebra and a brief introduction to tensor operations on vector spaces are presented.
	
	The idea of a group representation is introduced in section 3 as a homomorphism (a map respecting the group operation) bringing the elements of a group into the set of invertible transformations of some vector space along with a few basic consequences of the definition. Some typical representations follow with examples applied on some small cyclic and symmetric groups. The notion of subrepresentation, analogous to vector subspaces, is introduced along with tensor operations on representations, allowing us to construct new representations out of already known ones. Some representations are indivisible however and are called irreducible, inviting the main theorem of this text that any representation of a finite group in a finite dimensional vector space is composed of a finite number of constituent subrepresentations (Maschke's Theorem~\ref{thm:maschkes}), which along with Schur's Lemma (Theorem~\ref{thm:schur})  allows us to claim that these decompositions are unique.
	
	The last section is on character theory which strips back the ``big'' concepts of linear algebra and matrix calculations into the trace of those matrices, a metric related to the eigenvalues of a matrix, releasing us from the ambiguities of choices of bases. The typical representations found earlier are reevaluated and we see that the results allows us to find every irreducible representation of a group by providing a limit for the number of irreducibles, and also completely decompose any arbitrary representation into irreducibles with some ``character calculus''. The cyclic groups are quickly studied, and finally the symmetric groups of degree 3 and 4 are closely studied.

%% your first section
\clearpage{\thispagestyle{empty}\cleardoublepage}
\section{Introduction}

Lorem ipsum dolor sit amet, consectetur adipiscing elit, sed do eiusmod tempor incididunt ut labore et dolore magna aliqua.
Ut enim ad minim veniam, quis nostrud exercitation ullamco laboris nisi ut aliquip ex ea commodo consequat. 
Duis aute irure dolor in reprehenderit in voluptate velit esse cillum dolore eu fugiat nulla pariatur. 
Excepteur sint occaecat cupidatat non proident, sunt in culpa qui officia deserunt mollit anim id est laborum.

You can also use swedish letters ä ö å, Ä, Ö, Å. Unicode characters in general should work.


%Another section
\clearpage{\thispagestyle{empty}\cleardoublepage}
\section{Theorems and definitions}

\begin{definition}[Test av definition]
	Test av definition
\end{definition}

\begin{problem}[Test av problem]
	Test av problem
\end{problem}

\begin{conjecture}[Test av conjecture]
	Test av conjecture
\end{conjecture}

\begin{theorem}[Test av theorem]
	Test av theorem
\end{theorem}

\begin{proposition}[Test av proposition]
	Test av proposition
\end{proposition}

\begin{corollary}[Test av corollary]
	Test av corollary
\end{corollary}

\begin{lemma}[Test av lemma]
	Test av lemma
\end{lemma}

\begin{question}[Test av question]
	Test av question
\end{question}

\begin{remark}[Test av remark]
	Test av remark
\end{remark}

\begin{example}[Test av example]
	Test av example
\end{example}

\begin{notation}[Test av notation]
	Test av notation
\end{notation}

\begin{note}[Notering]
	This note's for you
\end{note}


\begin{definition}
	This is a definition!
\end{definition}

\begin{theorem}\label{thm: label}
	This is a theorem!
\end{theorem}

To cross reference: By Theorem~\ref{thm: label} we can say...
\begin{lemma}
	This is a lemma.
\end{lemma}
\begin{proof}
	Put here the proof of the lemma!
\end{proof}


\begin{proof}[Proof of Theorem \ref{thm: label}]
	Another proof
\end{proof}


\begin{corollary}
	here is a corollary!
\end{corollary}

If you want to write a beautiful equation in line do like that $e^{i\pi}+1=0$.
You might also want to write it on its own line
\[
e^{i\pi}+1=0.
\]
If you think that you are going to need to cross reference your equation:
\begin{equation}\label{eq:beautiful}
	e^{i\pi}+1=0
\end{equation}
and this is how you cross reference it \eqref{eq:beautiful}
You might also want to have aligned equations on more lines.
\begin{align}
	e^{i\pi}+1&=0\\
	(\cos\theta)^2+(\sin\theta)^2 &=1
\end{align}

If you need to cross reference something in the bibliography, use the cite command: \cite[p. 45]{Biggs2002}. 
The actual references are stored in a different file, called bibliographyFile.bib 

You can reference most things you found useful,
for example, books \cite{Biggs2002,Knuth1998ArtOfProgramming,LascouxPolynomialsBook}.
Research articles and preprints, \cite{AlexanderssonSulzgruber2019,Sherman2015x}.
Or another thesis, \cite{Teff2013}.

You can also cite online sources. For example, the Encyclopedia of integer sequences, \cite{OEIS},
or a YouTube video \cite{Bazett2018yt}. 
It is even possible to refer to a discussion you had with someone, \cite{PrivateCommunicationWithJustin}.
This is usually rare, it is better to thank them in the acknowledgements section (if they were helpful).



%another section
\clearpage{\thispagestyle{empty}\cleardoublepage}
\section{Tables and graphs}
Here a few examples of tables and graphs.
%A subsection
\subsection{Tables}
\begin{center}
	\begin{tabular}{c c c c c c c c}
		\arrayrulecolor{Azzurro}
		\hline
		{\bfseries Code} & {\bfseries CdL} & {\bfseries xxx} & {\bfseries $T_{setup/lotto}$} & {\bfseries $T_{lav/st}$} & {\bfseries $T_{proc/pezzo}$} & {\bfseries Quantitity} & {\bfseries $T_{tot}$}\\
		\hline
		100 & 4 & 250 & 25 & 0,5 & 0,6 & 1 & 0,6\\
		111 & 2 & 250 & 20 & 2 & 2,08 & 1 & 2,08 \\
		111 & 3 & 250 & 15 & 1,5 & 1,56 & 1 & 1,56 \\
		112 & 2 & 250 & 20 & 2,5 & 2,58 & 1 & 2,58 \\
		112 & 3 & 250 & 15 & 2 & 2,06 & 1 & 2,06\\
		113 & 3 & 500 & 15 & 1 & 1,03 & 2 & 2,06\\
		120 & 1 & 50 & 30 & 2 & 2,6 & 0,1 & 0,26\\
		121 & 1 & 25 & 30 & 3 & 4,2 & 0,1 & 0,42 \\
		121 & 1 & 25 & 30 & 2,5 & 3,7 & 0,1 & 0,37 \\
		\hline
	\end{tabular}
\end{center}



%another section
\clearpage{\thispagestyle{empty}\cleardoublepage}

\section{Insert a figure}
\begin{figure}[H]
	\centering
	\includegraphics[width=1\textwidth]{grafo.png}
	\caption{\label{fig:1} put the caption here}
\end{figure}

\subsection{Footnote}

You can create a footnote like this.\footnote{I created a footnote.}



\clearpage{\thispagestyle{empty}\cleardoublepage}
\section{Conclusion}

Remember the internet is full of places that can help you write a beautiful document with latex.
Some example are \url{overleaf.com} and \verb!stackexchange! just ask Google! 
In addition with an SU e-mail address you should be able to 
create (for free) a premium account on \url{overleaf.com}, 
so you do not need to download a tex editor and you will have
the full latex distribution at your hand (if you are connected to the internet).


$$ \NN \ \ZZ \ \QQ \ \RR \ \CC \ \KK \ \FF$$
$$ \1 \ \0 $$

$$ \GL(V),  \SL_n(\KK), \SymP^k V, \AltP^k V, \Hom(V,W), \SO, \SU, \id, \Stab(g), \Cent(g), \Norm(g),$$ $$ \Type(\sigma), \Dih_{2n}, \Sym_n, \Alt_n, \Klein, \Cyc_n, \sgn(\sigma), \Span(\bas_i)_i $$

$$ \bas \ihat \jhat \ghat \hhat \vhat \what \uhat \vvec \uvec \wvec $$
\clearpage{\thispagestyle{empty}}
\section{Preliminary topics}

\subsection{Linear algebra topics}

	\paragraph{Definition of kernel and image of a map}
	
	\begin{definition}[Kernel and image of a linear map]\label{def:kernelimage}
		Let $V$ and $W$ be two vector spaces and let $\varphi: V \rightarrow W$ be a linear map. Then the kernel and the image of the map are defined thusly:
		\begin{itemize}
			\item[i)] $\ker \varphi = \left\lbrace \vvec \in V \ \middle\vert \ \varphi(\vvec) = \0 \right\rbrace = \varphi^{-1}(\0)$.
			\item[ii)] $\im \varphi = \left\lbrace \wvec \in W \ \middle\vert \ \exists \vvec \in V \text{ s.t. } \varphi(\vvec) = \wvec \right\rbrace = \varphi(V)$.
		\end{itemize}
	\end{definition}
	
	\begin{remark}
		The kernel and the image of a linear map are subspaces of the domain and codomain of the map respectively, ie. $\ker \varphi$ is a subspace of $V$ and $\im \varphi$ is a subspace of $W$.
	\end{remark}
	
	\paragraph{Existence of complementary subspaces}
	
	\begin{corollary}\label{thm:compsubspaces}
		Let $V$ be a vector space and $W$ be a vector subspace of $V$. Then there exists a complementary vector subspace $W'$ in $V$ such that $W \cap W' = \emptyset$ and $W \cup W' = V$. This is equivalent to saying that $V$ is the direct sum of $W$ and $W'$, denoted as $V = W \oplus W'$.
	\end{corollary}
	\begin{proof}
		La la la
	\end{proof}
	
\subsection{Group theoretical topics}

	\subsubsection{Group action}
	
		We could let the group $G$ act on a set $X$ and study the action of $G$ on $X$ defined by
		\begin{align*}
			G \times X \rightarrow X
		\end{align*}
		and thereby permuting the elements of $X$ by for any $g \in G$,
		\begin{align*}
			(g,x) \mapsto g \cdot x = gx,
		\end{align*}
		for any $x \in X$.
	
	\subsubsection{The symmetric group}
	
		Recall that two elements $\sigma, \tau \in \Sym_n$ are conjugate if there exists a $\pi \in \Sym_n$ such that $\sigma = \pi \tau \pi^{-1}$. The conjugacy classes partition $\Sym_n$ into disjoint subsets. 
		
		\begin{proposition}[Conjugacy in $\Sym_n$]
			Any elements $\sigma, \tau \in \Sym_n$ are conjugate if and only if they are of the same cycle type.
		\end{proposition}
		
		\begin{definition}[Sign of a permutation]
			\begin{align*}
				\sgn(\sigma) = (-1)^k,
			\end{align*}
			where $k$ is the number of transpositions required to compose $\sigma$.
		\end{definition}
		
		\begin{proposition}\label{prop:signwelldefined}
			The sign of a permutation is well-defined. That is a permutation is either \textit{even} with sign +1 or \textit{odd} with sign -1.
		\end{proposition}
		\begin{proof}
			\cite[Thm.12.6.1.]{Biggs}
		\end{proof}
		
		
		

\clearpage{\thispagestyle{empty}}
\section{Representations of finite groups}

\marginnote{Include note on source!} 

\paragraph{First definitions.}

The main purpose of this text is to study linear representations of finite groups.

Let $G$ be a finite group written multiplicatively. Denote its order by $|G|$ and the action of its elements with $\cdot$, eg. $g \cdot x$.

Let $V$ be a finite-dimensional vector space over the field of complex numbers $\CC$. Its group of automorphisms is denoted by $\GL(V)$. If $V$ is provided with a basis, usually denoted $(\bas_i)_{i=1}^n$, where $n = \dim V$, then $\GL(V)$ is identified with the set of linear transformations of $V$, denoted $\GL_n(\CC)$\cite[18.1]{DummitFoote}.

\begin{definition}[Representation]\cite{Serre}
	Given a finite group $G$, a representation of $G$ in the vector space $V$ is a homomorphism
	\begin{align}\label{eq:DefRepr}
		\rho: G \rightarrow \GL(V).
	\end{align}
\end{definition}
For every $g$ in $G$, there is a linear map $\rho_g$ in $\GL(V)$ with an action on the elements of $V$ defined by 
\begin{align*}
	g \cdot \vvec = \rho_g (\vvec).
\end{align*}
The linearity of the map means that for any $\vvec, \wvec \in V$,
\begin{align}\label{eq:linearityOfg}
	g \cdot (a\vvec + b\wvec) = ag\cdot\vvec + bg\cdot\wvec,
\end{align} 
and also since it is a homomorphism, for any $g,h \in G$ we have
\begin{align}\label{eq:homomorphism}
	\rho_{gh} = \rho_g \rho_h.
\end{align}
From the homomorphism of $\rho$, two consequences follow.
\begin{proposition}
	For the identity element $e$ of $G$ and an arbitrary element $g$ with inverse $g^{-1}$ in $G$ we have,
	\begin{align}\label{eq:rhoIdAndInv}
		\rho_e = \id \quad \text{and} \quad \rho_g^{-1} = \rho_{g^{-1}},
	\end{align}
	where $\id$ is the identity transformation and $\rho_{g^{-1}}$ is the inverse of the representation of $g$.
\end{proposition}
\begin{proof}
	Take $g$ as an arbitrary element of $G$. The first identity follows from taking $h=e$ in Equation~\ref{eq:homomorphism}:
	\begin{align*}
		\rho_{eg} = \rho_e \rho_g.
		%\rho(eg) = \rho(e)\rho(g). %	\rho(g) \overset{!}{=} \rho(eg) = \rho(e)\rho(g)
	\end{align*}
	Since $eg=g$ for any $g\in G$ we must have that $\rho_g = \rho_e\rho_g$ which is true if and only if $\rho_e = \id$. Instead taking $h=g^{-1}$ we have:
	\begin{align*}
		\rho_{gg^{-1}} = \rho_g \rho_{g^{-1}},
		%\rho(g) \rho(g^{-1}) =  \rho(gg^{-1}~). %\overset{!}{=} \rho(e) = \id.
	\end{align*}
	but $gg^{-1} = e$ and $\rho_e = \id$, implying $\rho_{g^{-1}} = \rho_g^{-1}$.
\end{proof}

\begin{note}
	A vector space $V$ together with a group action fulfilling Equations \ref{eq:DefRepr}, \ref{eq:homomorphism} and \ref{eq:rhoIdAndInv} gives $V$ the structure of a $G$-module \cite[1.3]{Sagan}. To avoid a treatment on modules, if a vector space $V$ is a $G$-module it will be referred to as a \textit{representation space} of $G$.
\end{note}

\begin{note}
	The homomorphism $\rho$, the set of maps $\{\rho_g\}_{g \in G}$ and the representation space $V$ are interchangeably and abusively called the \textit{representation of $G$}.
\end{note}

\paragraph{Choice of basis.} 

If a basis $\{\bas_i\}_{i=1}^n$ is provided for a representation space, then a \textit{matrix representation} can be provided. In this case, $\rho_g$ is the atrix in $\GL_n(\CC)$ associated with the action of $g$ expressed in the basis provided. A matrix representation is not canonical and is dependent on the basis chosen.

\paragraph{Degree of a representation.} 

\begin{definition}[Degree of a representation]
	Let $V$ be a representation space of $G$. Then the dimension of $V$ is referred to as the degree of the representation.
\end{definition}	

\subsection{Tensor operations on representations}\label{sect:tensorrepr}

If we have two or more representations, we may construct additional ones using tensor operations. Given two representations $\rho^V: G \rightarrow \GL(V)$ and $\rho^W: G \rightarrow \GL(W)$ of a group $G$, the following are also representations.

\textbf{The direct sum of $V \oplus W$}, given by 
\begin{align*}
	\rho_{V \oplus W}(\vvec \oplus \wvec) = (\rho_V \oplus \rho_W)(\vvec \oplus \wvec) = \rho_V(\vvec) \oplus \rho_W(\wvec).
\end{align*}

By recursion, for a positive integer $n$, the \textbf{direct sum of $n$ copies of $V$}, denoted
\begin{align*}
	nV := \bigoplus_{i=1}^n V = \underset{\text{$n$ times}}{\underbrace{V \oplus \dots \oplus V}}.
\end{align*}

The \textbf{tensor product $V \otimes W$}, given by
\begin{align*}
	\rho_{V \otimes W}(\vvec \otimes \wvec) = (\rho_V \otimes \rho_W)(\vvec \otimes \wvec) = \rho_V(\vvec) \otimes \rho_W(\wvec).
\end{align*}

By recursion, for a positive integer $n$, the \textbf{$n$th tensor power of $V$}, denoted
\begin{align*}
	V^{\otimes n} := \bigotimes_{i=1}^n V = \underset{\text{$n$ times}}{\underbrace{V \otimes \dots \otimes V}}.
\end{align*}





%If we have two or more representations, we may construct additional ones using the tools of  Section~\ref{sect:tensoralgebra}. Given two representations $\rho^V: G \rightarrow \GL(V)$ and $\rho^W: G \rightarrow \GL(W)$ of a group $G$, the following are also representations.

%\begin{itemize}
%	\item The \emph{direct sum of $V \oplus W$} is a representation given by 
%	\begin{align*}
%		\rho_{V \oplus W}(\vvec \oplus \wvec) = (\rho_V \oplus \rho_W)(\vvec \oplus \wvec) = \rho_V(\vvec) \oplus \rho_W(\wvec).
%	\end{align*}
	
%	\item By recursion, for a positive integer $n$, the \emph{direct sum of $n$ copies of $V$}, denoted
%	\begin{align*}
%		nV := \bigoplus_{i=1}^n V = \underset{\text{$n$ times}}{\underbrace{V \oplus \dots \oplus V}},
%	\end{align*}
%	is a representation.
	
%	\item The \emph{tensor product $V \otimes W$} is a representation given by
%	\begin{align*}
%		\rho_{V \otimes W}(\vvec \otimes \wvec) = (\rho_V \otimes \rho_W)(\vvec \otimes \wvec) = \rho_V(\vvec) \otimes \rho_W(\wvec).
%	\end{align*}
	
%	\item By recursion, for a positive integer $n$, the \emph{$n$th tensor power of $V$}, denoted
%	\begin{align*}
%		V^{\otimes n} := \bigotimes_{i=1}^n V = \underset{\text{$n$ times}}{\underbrace{V \otimes \dots \otimes V}},
%	\end{align*}
%	is a representation. The zeroth power is by definition the ground field $\CC$ and the first power is $V$ itself.
%	\item The $n$th tensor power has two subrepresentations, the symmetric powers $\SymP^n V$ and the exterior powers $\AltP^n V$. \textit{Needs further explanation, especially the squares.}
%	\item After fixing a basis $(\bas_i)_{i=1}^k$ for $V$, a \emph{dual representation}, denoted $V^*$, can be defined by the dual space, spanned by the ...
	
%	\item Set of homomorphisms...
%\end{itemize}
\subsection{Examples}\label{sect:basicreprs}

This section follows \cite[1.2]{Serre}.

\subsubsection{Trivial representations}

%	A representation $\rho$ where every $g$ is such that $|\rho_g| = 1$, is called a unital representation.

\begin{example}[Trivial representation]\label{example:trivrepr}
	For any group $G$ there is a trivial degree 1 representation defined by mapping each element of $G$ to 1, ie.
	\begin{align*}
		\rho_g = 1,
	\end{align*}
	for every $g \in G$. It is clearly a homomorphism since for any $g,h \in G$ we have that
	\begin{align*}
		\rho_{gh} = 1 = 1 \cdot 1 = \rho_g \rho_h.
	\end{align*}
	This could be extended to any vector space by mapping $g$ to the identity transformation of that vector space\footnote{For a vector space of dimension $n$, the representation taking every $g$ in a group to the $n \times n$ identity matrix can be described, in the language of Section~\ref{sect:tensorrepr}, as the direct sum of $n$ copies of the trivial representation.}.
\end{example}

\begin{note}
	The trivial representation (the mapping is trivial) is not to be confused with the trivial space (the vector space containing only the zero vector).
\end{note}

\begin{example}[Alternating representation of $\Sym_n$]\label{example:altrepr}
	Choosing $G = \Sym_n$, another degree 1 representation can be found by studying the signs, or parities, of the elements of $\Sym_n$. By \cite[Thm.12.6.1.]{Biggs}, the sign of a permutation is well-defined, so for any two permutations $\sigma$ and $\tau$ with respective signs $(-1)^s$ and $(-1)^t$, their composition has the sign
	\begin{align*}
		\sgn(\sigma\tau) = (-1)^{s+t} = (-1)^s \cdot (-1)^t = \sgn (\sigma) \cdot \sgn (\tau),
	\end{align*}
	so clearly the map $\sgn: \Sym_n \rightarrow \{\pm 1\}$ is a homomorphism and thus a representation of degree 1, where even permutations are mapped to $+1$ and odd to $-1$.
\end{example}

\subsubsection{Degree 1 representations of $\Cyc_n$}\label{sect:deg1cycn}

Choose $G = \Cyc_n$, and let $g$ be a generator of $\Cyc_n$ such that 
\begin{align*}
	\Cyc_n = \{ e, g, g^2, \dots, g^{n-1}\}
\end{align*}
and $g^n = e$. Consider a map $\rho: \Cyc_n \rightarrow \CC$ defined as a homomorphism by $\rho_{g^a}\rho_{g^b} = \rho_{g^{a+b}}$ for some integers $a,b$, then by Equation~\ref{eq:rhoIdAndInv} we have that $\rho_e = 1$, but $g^n = e$ and by induction on Equation~\ref{eq:homomorphism} we have that $\rho_{g^n} = (\rho_g)^n$, so then we must have that $(\rho_g)^n = 1$, that is $\rho_g$ is mapped to a $n$th root of unity.
In conclusion, for $\Cyc_n$ we have found $n$ representations of degree 1, each mapping $g$ to a $n$th root of unity and the powers of $g$ to the corresponding powers of the root of unity. 

\begin{example}[$\Cyc_3$]
	The three third roots of unity are $1$, $\omega = \frac{-1+i\sqrt{3}}{2}$ and $w^2 = \frac{-1-i\sqrt{3}}{2}$, so the three representations of $\Cyc_3$ are presented in table \ref{table:Cyc3}. Note that $\rho_1^* = \rho_2$.
	
%	\begin{table}[hbt!]
%		\begin{tabular}{c | c c c}
%			$\Cyc_3$ & $e$ & $g$        & $g^2$      \\ \hline
%			$\rho_0$          & 1   & 1          & 1          \\
%			$\rho_1$          & 1   & $\omega$   & $\omega^2$ \\
%			$\rho_2$          & 1   & $\omega^2$ & $\omega$
%		\end{tabular}
%		\centering
%		\caption{Three representations of $\Cyc_3$. $\omega = e^{2 \pi i/3}$}
%		\label{table:Cyc3}
%	\end{table}
\end{example}

\begin{example}[$\Cyc_4$]
	The four fourth roots of unity are $1,i,-1$ and $-i$ and the four corresponding representations of $\Cyc_4$ are presented in Table~\ref{tbl:cyc4}. Note that $\rho_1^* = \rho_3$.
%	\begin{table}[hbt!]
%		\centering
%		\begin{tabular}{c | c c cl}
%			$\Cyc_4$ & $e$ & $g$  & $g^2$ & $g^3$ \\ \hline
%			$\rho_0$           & 1   & 1    & 1     & 1     \\
%			$\rho_1$           & 1   & $i$  & $-1$  & $-i$  \\
%			$\rho_2$           & 1   & $-1$ & $1$  & $-1$   \\
%			$\rho_3$           & 1   & $-i$ & $-1$   & $i$
%		\end{tabular}
%		\caption{Four representations of $\Cyc_4$.}
%		\label{tbl:cyc4}
%	\end{table}
\end{example}
	
\begin{example}[$\Cyc_5$]
	The five fifth roots of unity are $\exp{\frac{2\pi im}{5}}$, $0 \leq m \leq 4$, so five representations of $\Cyc_5$ are presented in table \ref{table:Cyc5}. Note that $\rho_1^* = \rho_4$ and $\rho_2^* = \rho_3$.
	
%	\begin{table}[hbt!]
%		\begin{tabular}{c | c c c c c}
%			
%			$\Cyc_5$ & $e$ & $g$        & $g^2$      & $g^3$      & $g^4$      \\ \hline
%			$\rho_0$            & 1   & 1          & 1          & 1          & 1          \\
%			$\rho_1$            & 1   & $\omega^1$ & $\omega^2$ & $\omega^3$ & $\omega^4$ \\
%			$\rho_2$            & 1   & $\omega^2$ & $\omega^4$ & $\omega^1$ & $\omega^3$ \\
%			$\rho_3$            & 1   & $\omega^3$ & $\omega^1$ & $\omega^4$ & $\omega^2$ \\
%			$\rho_4$            & 1   & $\omega^4$ & $\omega^3$ & $\omega^2$ & $\omega^1$
%		\end{tabular}
%		\centering
%		\caption{Five representations of $\Cyc_5$. $\omega = e^{2 \pi i/5}$.}
%		\label{table:Cyc5}
%	\end{table}
\end{example}



\begin{table}[hbt!]
	\parbox[s]{.31\linewidth}{
		\centering
		\begin{tabular}{c | c c c}
			$\Cyc_3$ & $e$ & $g$        & $g^2$      \\ \hline
			$\rho_0$          & 1   & 1          & 1          \\
			$\rho_1$          & 1   & $\omega$   & $\omega^2$ \\
			$\rho_2$          & 1   & $\omega^2$ & $\omega$
		\end{tabular}
		\vfill
		\caption{Three reprs. of $\Cyc_3$. $\omega = e^{2 \pi i/3}$.}
		\label{table:Cyc3}
	}
	\hfill
	\parbox[s]{.31\linewidth}{
		\centering
		\begin{tabular}{c | c c cl}
			$\Cyc_4$ & $e$ & $g$  & $g^2$ & $g^3$ \\ \hline
			$\rho_0$           & 1   & 1    & 1     & 1     \\
			$\rho_1$           & 1   & $i$  & $-1$  & $-i$  \\
			$\rho_2$           & 1   & $-1$ & $1$  & $-1$   \\
			$\rho_3$           & 1   & $-i$ & $-1$   & $i$
		\end{tabular}
		\caption{Four reprs. of $\Cyc_4$.}
		\label{tbl:cyc4}
	}
	\hfill
	\parbox[s]{.31\linewidth}{
		\centering
		\begin{tabular}{c | c c c c c}
			$\Cyc_5$ & $e$ & $g$        & $g^2$      & $g^3$      & $g^4$      \\ \hline
			$\rho_0$            & 1   & 1          & 1          & 1          & 1          \\
			$\rho_1$            & 1   & $\omega$ & $\omega^2$ & $\omega^3$ & $\omega^4$ \\
			$\rho_2$            & 1   & $\omega^2$ & $\omega^4$ & $\omega$ & $\omega^3$ \\
			$\rho_3$            & 1   & $\omega^3$ & $\omega$ & $\omega^4$ & $\omega^2$ \\
			$\rho_4$            & 1   & $\omega^4$ & $\omega^3$ & $\omega^2$ & $\omega$
		\end{tabular}
		\caption{Five reprs. of $\Cyc_5$. $\omega = e^{2 \pi i/5}$.}
		\label{table:Cyc5}
	}
\end{table}



\subsubsection{Permutation representation}

	Given a group $G$, we chose a set $X$ that $G$ acts on by permutation. Let $V$ be a vector space spanned by a natural basis $(\bas_x)_{x \in X}$, then we have a representation
	\begin{align*}
		\rho:G \rightarrow \GL(V)
	\end{align*} 
	defined by its action on the basis vectors by, for any $g \in G$, 
	\begin{align*}
		\rho_g: \bas_x \mapsto \bas_{gx},
	\end{align*}
	that is $\rho$ inherited the group action of $G$ on $X$. It is a homomorphism since for any $g,h \in G$ we have
	\begin{align*}
		g \cdot (h \cdot \bas_x) &= g \cdot \bas_{hx} \\
		&= \bas_{ghx} \\
		&= (gh) \cdot \bas_x
	\end{align*}
	for any $x \in X$.%\marginnote{We have imposed a $G$-module structure on $V$.}
	In the $(\bas_x)_{x \in X}$-basis, $\rho_g$ are permutation matrices, which have one once in every row and column, and the rest of the entries are zero.
	
	Letting $G = \Sym_n$, it would be appropriate to choose the set $X = \{1, 2, \dots, n\}$ and to let any $\sigma \in \Sym_n$ permute any $1 \leq i \leq n$ by $i \mapsto \sigma(i)$. Choosing a basis $(\bas_i)_{i=1}^n$ to span a vector space $V$, we define a homomorphism
	\begin{align*}
		\rho: \Sym_n \rightarrow \GL(V)
	\end{align*}
	defined for any $\sigma \in \Sym_n$ as
	\begin{align*}
		\rho_\sigma: \bas_i \mapsto \bas_{\sigma(i)}
	\end{align*}
	for any $1 \leq i \leq n$. In this basis, the corresponding set of matrix representations are the permutation matrices
	\begin{align*}
		\rho_\sigma = (r_{ij})_{n \times n} \text{, where } r_{ij} = \delta_{j,\sigma(i)} = \begin{cases}
			1 \text{, if $j = \sigma(i)$,} \\
			0 \text{, otherwise,}
		\end{cases}
	\end{align*}
%	where $\delta_{k,l}$ is the Kronecker delta which vanishes if $k \neq l$ and is equal to one if $k=l$.
	in which the $i$th column has a 1 in the $\sigma(i)$th row, and the rest of the rows have a 0.
	
	Under this action, a vector
	\begin{align*}
		(a_1, a_2, \dots, a_n) = a_1 \bas_1 + a_2 \bas_2 + \dots + a_n \bas_n \in V
	\end{align*}
	is permuted to 
	\begin{align*}
		(a_{\sigma^{-1}(1)}, a_{\sigma^{-1}(2)}, \dots, a_{\sigma^{-1}(n)}) = a_1 \bas_{\sigma(1)} + a_2 \bas_{\sigma(2)} + \dots + a_n \bas_{\sigma(n)}  \in V.
	\end{align*}
	
%	\marginnote{Yet another interpretation is that $\sigma$ permutes the coordinate axes of a $n$-dimensional space.}
	
	The dimension of $V$ is $|X|$, for example the permutation representation of $\Sym_n$ is of degree $n$.
	
	\begin{example}[Permutation representation of $\Sym_2$]
		The symmetric group of degree 2 has two elements, $\Sym_2 = \{(1), (1,2)\}$ and their matrix representations in $\bas_1,\bas_2$-space are presented in Table~\ref{table:permS2}.
		\begin{table}[hbt!]
			\centering
			\begin{tabular}{c c}
				$\rho_{(1)} =
				\begin{pmatrix}
					1 & 0 \\ 0 & 1
				\end{pmatrix}$, &
				$\rho_{(1,2)} =
				\begin{pmatrix}
					0 & 1 \\ 1 & 0
				\end{pmatrix}$
			\end{tabular}
			\caption{Matrix representations of $\Sym_2$}
			\label{table:permS2}
		\end{table}
	\end{example}
	
	\begin{example}[Permutation representation of $\Sym_3$]\label{ex:permS3}
		Likewise, representations of $\Sym_3$ are presented in Table~\ref{table:permS3}.
		\begin{table}[hbt!]
			\centering
			\begin{tabular}{r r r}
				$\rho_{(1)} = 
				\begin{pmatrix}
					1 & 0 & 0 \\
					0 & 1 & 0 \\
					0 & 0 & 1
				\end{pmatrix}$, & 
				$\rho_{(1,2,3)} = 
				\begin{pmatrix}
					0 & 0 & 1 \\
					1 & 0 & 0 \\
					0 & 1 & 0
				\end{pmatrix}$, & 
				$\rho_{(1,3,2)} = 
				\begin{pmatrix}
					0 & 1 & 0 \\
					0 & 0 & 1 \\
					1 & 0 & 0
				\end{pmatrix}$, \\ & & \\
				$\rho_{(1,2)} = 
				\begin{pmatrix}
					0 & 1 & 0 \\
					1 & 0 & 0 \\
					0 & 0 & 1
				\end{pmatrix}$, &
				$\rho_{(1,3)} = 
				\begin{pmatrix}
					0 & 0 & 1 \\
					0 & 1 & 0 \\
					1 & 0 & 0
				\end{pmatrix}$, &
				$\rho_{(2,3)} = 
				\begin{pmatrix}
					1 & 0 & 0 \\
					0 & 0 & 1 \\
					0 & 1 & 0
				\end{pmatrix}$.
			\end{tabular}
			\caption{Matrix representations of $\Sym_3$}
			\label{table:permS3}
		\end{table}
	\end{example}
	
	\begin{example}[Permutation representation of $\Sym_4$]
		Yet again, matrix representations of some elements of $\Sym_4$ are presented in Table~\ref{table:permS4}. The elements chosen are one representative from every conjugacy class of $\Sym_4$.
		\begin{table}[hbt!]
			\centering
			\begin{tabular}{r r r}
				$\rho_{(1)} = \left(\begin{smallmatrix}
					1 & 0 & 0 & 0 \\
					0 & 1 & 0 & 0 \\
					0 & 0 & 1 & 0 \\
					0 & 0 & 0 & 1
				\end{smallmatrix}\right)$,  &
				$\rho_{(1,2)} = \left(\begin{smallmatrix}
					0 & 1 & 0 & 0 \\
					1 & 0 & 0 & 0 \\
					0 & 0 & 1 & 0 \\
					0 & 0 & 0 & 1
				\end{smallmatrix}\right)$, &
				$\rho_{(1,2)(3,4)} = \left(\begin{smallmatrix}
					0 & 1 & 0 & 0 \\
					1 & 0 & 0 & 0 \\
					0 & 0 & 0 & 1 \\
					0 & 0 & 1 & 0
				\end{smallmatrix}\right)$, \\ & & \\ &
				$\rho_{(1,2,3)} = \left(\begin{smallmatrix}
					0 & 0 & 1 & 0 \\
					1 & 0 & 0 & 0 \\
					0 & 1 & 0 & 0 \\
					0 & 0 & 0 & 1
				\end{smallmatrix}\right)$, &
				$\rho_{(1,2,3,4)} = \left(\begin{smallmatrix}
					0 & 0 & 0 & 1 \\
					1 & 0 & 0 & 0 \\
					0 & 1 & 0 & 0 \\
					0 & 0 & 1 & 0
				\end{smallmatrix}\right)$.
			\end{tabular}
			\caption{Some matrix representations of $\Sym_4$}
			\label{table:permS4}
		\end{table}
	\end{example}

\subsubsection{Regular representation}

	Following the same reasoning as in the previous section, but instead we let $G$ act on itself, ie. $X = G$. The corresponding vector space $V$ is spanned by the basis $(\bas_g)_{g \in G}$ constructed from the elements of $G$. A representation of $G$ in $V$ is then a map
	\begin{align*}
		\rho: G \rightarrow \GL(V)
	\end{align*}
	defined by
	\begin{align*}
		g \cdot \bas_h = \bas_{gh}
	\end{align*}
	for any $g,h \in G$.%\marginnote{We have created a Group Algebra $\CC[G]$.}
	
	The dimension of $V$ is $|G|$, for example the regular representation of $\Cyc_n$ is of degree $n$ and for $\Sym_n$ it is of degree $n!$, a number which grows increasingly quick for larger $n$, however the regular representation will be shown to be important later.
	
	\begin{example}[Regular representation of $\Cyc_3$]\label{example:regCyc3}
		For $G = \Cyc_3$ we calculate the action of $\Cyc_3$ on the $\bas, \ghat, \ghat_2$-basis thusly:
		\begin{align*}
			\begin{cases}
				e \cdot \bas = \bas, \\
				e \cdot \ghat = \ghat, \\
				e \cdot \ghat^2 = \ghat^2,
			\end{cases} \quad \begin{cases}
				g \cdot \bas = \ghat, \\
				g \cdot \ghat = \ghat^2, \\
				g \cdot \ghat^2 = \bas,
			\end{cases} \quad \begin{cases}
			g^2 \cdot \bas = \ghat^2, \\
			g^2 \cdot \ghat = \bas, \\
			g^2 \cdot \ghat^2 = \ghat,
			\end{cases}
		\end{align*}
		hence the permutation matrices in the $\bas, \ghat, \ghat_2$-basis are
		\begin{align*}
			\rho_e = \id, \quad \rho_g = \begin{pmatrix}
				0&0&1 \\ 1&0&0 \\ 0&1&0
			\end{pmatrix}, \quad \rho_{g^2} = \begin{pmatrix}
			0&1&0 \\ 0&0&1 \\ 1&0&0 
			\end{pmatrix}.
		\end{align*}
		Let's choose a new basis inspired by the third roots of unity 1, $\omega$ and $\omega^2$:
		\begin{align*}
			\begin{cases}
				\fbas_1 = \bas + \ghat + \ghat^2, \\
				\fbas_2 = \bas + \omega^2\ghat + \omega\ghat^2, \\
				\fbas_3 = \bas + \omega\ghat + \omega^2\ghat^2,
			\end{cases}
		\end{align*}
		where $\omega = e^{2 \pi i/3}$, corresponding to the change-of-basis matrix 
		\begin{align*}
			P = \begin{pmatrix}
				1&1&1 \\
				1&\omega^2&\omega \\
				1&\omega&\omega^2
			\end{pmatrix}, \quad \text{with inverse} \quad P^{-1} = \frac{1}{3}\begin{pmatrix}
				1&1&1 \\
				1&\omega&\omega^2 \\
				1&\omega^2&\omega
			\end{pmatrix}.
		\end{align*}
		Now, the matrices of $\Cyc_3$ in this new basis are
		\begin{align*}
			\rho_e &= P^{-1} \cdot \id \cdot P = \id, \\ 
			\rho_g &= P^{-1}\cdot\begin{pmatrix}
				0&0&1 \\ 1&0&0 \\ 0&1&0
			\end{pmatrix}\cdot P = \begin{pmatrix}
			1&0&0 \\ 0&\omega&0 \\ 0&0&\omega^2
			\end{pmatrix}, \\
			\rho_{g^2} &= P^{-1}\cdot\begin{pmatrix}
				0&1&0 \\ 0&0&1 \\ 1&0&0 
			\end{pmatrix}\cdot P = \begin{pmatrix}
			1&0&0 \\ 0&\omega^2&0 \\ 0&0&\omega
			\end{pmatrix}.
		\end{align*}
		We can clearly see that the regular representation of $\Cyc_3$ is the direct sum of the three degree 1 representations found in Table~\ref{table:Cyc3}, ie. $\rho_g = (1) \oplus (\omega) \oplus (\omega^2)$ after the change of basis.
	\end{example}
	
	\begin{example}[Regular representation of $\Sym_3$]\label{example:regSym3}
		Defining a homomorphism \begin{align*}
			\rho: S_3 \rightarrow \GL(V)
		\end{align*} on the vector space spanned by the basis\footnote{Explicitly, this is $\left\lbrace \bas_{(1)}, \bas_{(1,2)}, \bas_{(1,3)}, \bas_{(2,3)}, \bas_{(1,2,3)}, \bas_{(1,3,2)} \right\rbrace$.} $(\bas_{\sigma})_{\sigma \in \Sym_3}$ defined for every $\sigma$ by \begin{align*}
		\rho_\sigma \cdot \bas_\tau = \bas_{\sigma \circ \tau} = \bas_{\sigma\tau}
		\end{align*} for some $\tau \in \Sym_3$, we arrive at a degree 6 representation of $\Sym_3$. %It is clearly a larger degree representation of $S_n$ than any we have found so far.% Later we will find that for $\Sym_3$, $R$ is the direct sum of the trivial representation, the alternating representation and two copies of the standard representations, that is $S = T \oplus A \oplus 2S$.
		
		\textit{To be added: Table of $6 \times 6$-matrices, one per conjugacy class for example. Also: Show that perm.repr. of $\Sym_3$ is completely decomposed to trivial + alternating + two copies of standard, especially since the multiplicity of the standard representation is $>1$.}
	\end{example}
	
	

%\subsubsection{Coset representations}
%
%Let $H$ be a subgroup of $G$. Let $\mathcal{H} = \{g_1H, g_2H, \dots, g_kH\}$ be the set of left cosets of $H$ in $G$ and define the action of $G$ on $\mathcal{H}$ by $g(g_iH) = (gg_i)H$ for any $g\in G$ and $g_iH \in \mathcal{H}$. The $g_i$ are called transversals of their respective $g_iH$. Construct a vector space $V$ by spanning it with the transversals as a basis, denoted $(\ghat_i)_{i=1}^k$ and define the group action of $G$ on $V$ by replicating the group action of $G$ on $\mathcal{H}$.\cite[Example 1.3.5.]{Sagan}
%
%\begin{example}[Trivial representation again]
%	Choosing $H = G$ (the non-proper subgroup), the coset representation reduces to the trivial representation.
%\end{example}
%
%\begin{example}[Regular representation again]
%	Choosing $H = \{e\}$ (the trivial subgroup), then $\mathcal{H} = G$ and we construct a vector space $V$ using the elements of $G$ as a basis, hence we have found the regular representation.
%\end{example}
%
%\begin{example}[Choosing a smarter subgroup of $\Sym_3$]
%	As laid out in~\cite[Example 1.3.5.]{Sagan}, choosing the subgroup of $\Sym_3$ to be $H = \{(1), (2,3)\}$, then the cosets of $H$ are $\mathcal{H} = \{ H, (1,2)H, (1,3)H \}$, and computing the matrices will yield the permutation representation of $\Sym_3$.
%\end{example}
\subsection{Subrepresentations}\label{sect:subrepr}

As laid out in \cite[Sect.?.?.]{FultonHarris}, we are looking for representations which are said to be ``atomic'' and conversely, for any arbitrary representation we wish to find how it is composed of these indecomposable representations. To proceed we need the notion of a vector space map that respects or conserves the group action.	

\begin{definition}[$G$-linear map]\cite{FultonHarris}
	Let $V$ and $W$ be two representation spaces of a group $G$. A vector space map 
	\begin{align*}
		\varphi: V \rightarrow W
	\end{align*} 
	is called a $G$-linear map if it commutes with the group action of $G$, ie. for any $\vvec \in V$ and $g \in G$ we have
	\begin{align*}
		\varphi( g \cdot \vvec ) = g \cdot \varphi (\vvec),
	\end{align*}
	or in terms of the maps $\rho_g^V: G \rightarrow \GL(V)$ and $\rho_g^W: G \rightarrow \GL(W)$,
	\begin{align*}
		\varphi \circ \rho^V_g (\vvec) = \rho^W_g \circ \varphi (\vvec).
	\end{align*}
	Equivalently one can say that the diagram in Figure~\ref*{fig:commapmorp} is commutes for every $g \in G$.
	\begin{figure}[hbt!]
		\centering
%		\[\begin{tikzcd}
%			V \arrow{r}{\varphi} \arrow[swap]{d}{g \cdot} & W \arrow{d}{g \cdot} \\
%			V \arrow{r}{\varphi} & W
%		\end{tikzcd}\]
		% https://q.uiver.app/#q=WzAsNCxbMCwwLCJWIl0sWzIsMCwiVyJdLFswLDIsIlYiXSxbMiwyLCJXIl0sWzAsMiwiZ1xcY2RvdCIsMl0sWzEsMywiZ1xcY2RvdCJdLFsyLDMsIlxcdmFycGhpIiwyXSxbMCwxLCJcXHZhcnBoaSJdXQ==
		\[\begin{tikzcd}
			V && W \\
			\\
			V && W
			\arrow["\varphi", from=1-1, to=1-3]
			\arrow["g\cdot"', from=1-1, to=3-1]
			\arrow["g\cdot", from=1-3, to=3-3]
			\arrow["\varphi"', from=3-1, to=3-3]
		\end{tikzcd}\]
		\caption{The map $\varphi$ is $G$-linear if the diagram commutes.}
		\label{fig:commapmorp}
	\end{figure}
\end{definition}

Consider the case where $W$ is a $G$-invariant\footnote{For $W$ to be a $G$-invariant space, it means that for every $g \in G$, $g \cdot \wvec \in W$ for any $\wvec \in W$ or that $\rho_g(W) \subseteq W$ for any $g\in G$.} vector subspace of $V$% and let $\pi:V \rightarrow W$ be the projection map of $V$ onto $W$. 
and let $\rho$ be a representation of $G$ in $V$. Then the restriction of $\rho$ to $W$, here denoted $\rho|_W$, is a isomorphism of $W$ onto itself since $W$ is $G$-invariant,
%(the action of $g$ on any $\wvec \in W$ will return the same $\wvec$), 
and thus $\rho_{gh}|_W = \rho_g|_W \cdot \rho_h|_W$, hence $\rho|_W: G \rightarrow \GL(W)$ is a representation of $G$ in the subspace $W$ of $V$, motivating the following definition:

\begin{definition}[Subrepresentation]\label{def:subrepr}
	Let $G$ be a finite group and $\rho$ be a representation of $G$ in a vector space $V$. A restriction of $\rho$ to a $G$-invariant vector subspace $W$ of $V$ is called a subrepresentation of $\rho$.
\end{definition}

%In other words, a subrepresentation is a $G$-invariant vector subspace of a ``parent'' representation. 

\begin{example}[Trivial subspaces]
	Any representation has itself as well as the zero space as subrepresentations.
\end{example}

\subsubsection{Some proper and non-trivial subrepresentations}

Let $G$ be any group and let us study its permutation and regular representations. 

\begin{example}[Trivial representation inside the permutation representation]\cite[Example 1.4.3.]{Sagan}\label{example:trivpermrepr}
	Let $G$ act on a set $X = \{x_1, x_2, \dots x_k\}$, where $k = |X|$, and let $V$ be the vector space spanned by the basis $(\xhat_1, \xhat_2, \dots, \xhat_k)$. Consider the one-dimensional subspace of $V$ spanned by the sum of all basis vectors, ie. 
	\begin{align*}	
		W = \Span\{ \xhat_1 + \xhat_2 + \dots + \xhat_k \}.
	\end{align*}A vector $\wvec \in W$ is a scalar multiple of this sum, and for any $g \in G$, the action of $g$ on $\wvec$ will simply reorder this sum and return the same $\wvec$, that is for any $g \in G$, we have that $\rho_g|_W = 1$. Thus $W$ is a $G$-invariant subspace of $V$ and the permutation representation has the trivial representation as a subrepresentation.
\end{example}

\begin{example}[Trivial representation inside the regular representation]\cite[Example 1.4.4.]{Sagan}
	Similarly to the last example, we span a vector space $V$ by a basis $(\ghat_1, \ghat_2, \dots, \ghat_{k})$ for every $g_i \in G$, where $k = |G|$, and consider the subspace 
	\begin{align*}
		W = \Span\{ \ghat_1 + \ghat_2 + \dots + \ghat_k \}. 
	\end{align*}
	Completely analogously to the last example, $W$ is shown to be a $G$-invariant subspace of $V$ and the regular representation also has the trivial representation as a subrepresentation.
\end{example}

\begin{example}[Alternating representation inside the regular representation of $\Sym_n$]\cite[Example 1.4.4.]{Sagan}
	Let $G = \Sym_n$, then for the regular representation, $V$ is spanned by a basis vector for every $\sigma \in \Sym_n$. Let $W$ be the subspace of $V$ spanned by the sum $\sum_{\sigma \in \Sym_n} \sgn(\sigma) \bas_{\sigma}$. The action of a $\tau \in \Sym_n$ on a $\wvec \in W$ is
	\begin{align*}
		\tau (\wvec) &= \tau \cdot a\sum_{\sigma \in \Sym_n} \sgn(\sigma) \bas_{\sigma} \\
		&= a \sum_{\sigma \in \Sym_n} \sgn(\sigma) \bas_{\tau\sigma} \\
		&= a \frac{1}{\sgn\tau} \sum_{\sigma \in \Sym_n} \sgn(\tau\sigma) \bas_{\tau\sigma} 
	\end{align*}
	where $a \in \CC$. Let $\pi = \tau\sigma$, then the sum can instead be over all these $\pi \in \Sym_n$, hence $W$ is a $G$-invariant subspace of $V$ and the alternating representation is also found inside the regular representation of the symmetric group.
\end{example}

\subsubsection{Subrepresentations as a kernel}

From linear algebra, we are familiar with the kernel and image of a map (see Definition~\ref{def:kernelimage}).

\begin{proposition}[Kernel and image of a $G$-linear map]\label{prop:kernelimagelinearmap}
	Let $V$ and $W$ be representation spaces of a group $G$ and let $\varphi: V \rightarrow W$ be a $G$-linear map. Then the kernel and the image of $\varphi$ are also representations of $G$ and more specifically, $\ker \varphi$ is a subrepresentation of $V$ and $\im \varphi$ is a subrepresentation of $W$.
\end{proposition}
\begin{proof}
	\begin{itemize}
		\item[i)] Let $\vvec \in \ker \varphi$. Then $g \cdot \varphi (\vvec) = \0$ since by definition $\varphi(\vvec) = \0$, but then we must also have $\varphi ( g \cdot \vvec ) = \0$ since $\varphi$ is a $G$-linear map, implying that $g \cdot \vvec \in \ker \varphi$. Since this holds for any $g \in G$ and $\vvec \in \ker \varphi$, the kernel of $\varphi$ is a $G$-invariant subspace of $V$ and thus a subrepresentation of $V$. %\marginnote{Subrepr. not yet defined}
		\item[ii)] For a $\wvec \in \im \varphi$, by definition there exists a $\vvec \in V$ such that $\varphi(\vvec) = \wvec$. Since $g \cdot \vvec \in V$, then we have that $\varphi ( g \cdot \vvec) \in W$. This implies that $g \cdot \varphi (\vvec)$ is also in $W$. Since this holds for any $g \in G$ and $\vvec \in V$, the image of $\varphi$ is a $G$-invariant subspace of $W$ and a thus a subrepresentation of $W$.	\qedhere
		%			\item[i)] Let $\vvec \in \ker \varphi$. Then $\rho_g \circ \varphi (\vvec) = \0$, but then we must also have $\varphi \circ \rho_g (\vvec) = \0$, implying that $\rho_g (\vvec) \in \ker \varphi$. Since this holds for any $g \in G$ and any $\vvec \in \ker \varphi$, the kernel of $\varphi$ is a $G$-invariant subspace of $V$ and thus a subrepresentation of $V$.
		%			\item[ii)] For a $\wvec \in \im \varphi$, there exists a $\vvec \in V$ such that $\varphi(\vvec) = \wvec$. Since $\rho_g(\vvec) \in V$, then we have that $\varphi \circ \rho_g(\vvec) \in W$. This implies that $\rho_g \circ \varphi (\vvec)$ is also in $W$. Since this holds for any $g \in G$ and $\vvec \in V$, the image of $\varphi$ is a $G$-invariant subspace of $W$ and a thus a subrepresentation of $W$.	\qedhere
	\end{itemize}
\end{proof}

Recall from linear algebra the existence of complementary vector subspaces of a vector space. Is there a similar property of subrepresentations? Consider the case when $W$ is a subrepresentation of $V$ and let $\pi: V \rightarrow W$ be the projection of $V$ onto $W$. Recall (Corollary~\ref{thm:compsubspaces})~that for each subspace $W\leq V$ there exists a complementary and disjoint subspace $W' \leq V$ such that $V = W \oplus W'$. This means that every $\vvec \in V$ is partitioned into $\wvec \oplus \wvec'$ where $\wvec \in W$ and $\wvec' \in W'$ such that 
\begin{align*}
	\begin{cases}
		\pi(\vvec) = \wvec, \\
		\pi(\wvec) = \wvec \text{ and } \\
		\pi(\wvec') = \0.
	\end{cases}
\end{align*} 
The image of the projection is clearly $W$ and the kernel is $W'$. But is $\pi$ a $G$-linear map?
\begin{proposition}[Existence of complementary subrepresentations]\label{prop:complsubrepr}
	Let $\rho$ be a representation of a finite group $G$ in a finite-dimensional vector space $V$. Let $W$ be a $G$-invariant subspace of $V$. Then there exists a $G$-invariant subspace $W'$ of $V$ complementary to $W$ such that $V = W \oplus W'$.
\end{proposition}
\begin{proof}
	\textit{THIS PROOF MISSES SOME PARTS: WE NEED TO SHOW THAT $\barpi$ IS STILL A PROJECTION $V \rightarrow W$.}
	By Proposition~\ref{prop:kernelimagelinearmap} we know that a $G$-linear map from one representation space $V$ to another $W$ has its kernel and image as subrepresentations of $V$ and $W$ respectively. If we can find such a $G$-linear map from $V$ to a $G$-invariant subspace $W \leq V$ we are done.
	The projection $\pi$ may not generally be a $G$-linear map.% since $\pi(g \cdot \vvec)$ is not generally the same as $g \cdot \pi(\vvec)$. 
	\textit{NEEDS A COUNTEREXAMPLE?}
%	for a $\vvec \in V$ we have that \marginnote{There is a mistake here}
%	\begin{align*}
%		\pi(g \cdot \vvec) = \pi(\wvec) = \wvec,
%	\end{align*}
%	but 
%	\begin{align*}
%		g \cdot \pi(\vvec) = g \cdot \wvec,
%	\end{align*}
%	and the equality $g \cdot \wvec = \wvec$ is not generally true for every $g \in G$.
	Instead, consider taking the \textit{average} of $\pi$ over $G$,
	\begin{align*}
		\barpi = \frac{1}{|G|} \sum_{g \in G} \rho_g \cdot \pi \cdot \rho_g^{-1},
	\end{align*}
	and see if it conserves the action of $G$. That is, does $\barpi\cdot \rho_g = \rho_g \cdot \barpi$ hold for every $g \in G$? Equivalently, let's consider
	\begin{align*}
		\rho_g \cdot \barpi \cdot \rho_g^{-1} &= \frac{1}{|G|} \sum_{h \in G} \rho_g \cdot \rho_h \cdot \pi \cdot \rho_h^{-1} \cdot \rho_g^{-1} & \text{(Def. of $\barpi$)}\\
		&= \frac{1}{|G|} \sum_{h \in G} \rho_{gh} \cdot \pi \cdot \rho_{gh}^{-1} & \text{($\rho$ homom., $(gh)^{-1}=h^{-1}g^{-1}$)}\\
		&= \frac{1}{|G|} \sum_{g' \in G} \rho_{g'} \cdot \pi \cdot \rho_{g'}^{-1} & \text{(Let $g' = gh$)} \\
		&= \barpi.
	\end{align*}
	So, $\barpi \cdot \rho_g = \rho_g \cdot \barpi$ for every $g \in G$ and thus $\barpi: V \rightarrow W$ is a $G$-linear map, which means that its kernel, $\ker \barpi$, identified to be the complement $W'$ of $W$ in $V$, is also a $G$-invariant subspace of $V$ and thereby a subrepresentation of $V$.
\end{proof}

\subsubsection{The standard represention of $\Sym_n$}

In Example~\ref{example:trivpermrepr}, the trivial representation was found to be a subrepresentation of the permutation representation. Let $G= \Sym_n$ and let $V$ be the permutation representation and assign it the basis $(\bas_i)_{i=1}^n$. The one-dimensional subspace $W$ spanned by the sum of all $\bas_i$ was found to be the trivial representation, then by Proposition~\ref{prop:complsubrepr}, there exists another subrepresentation $W'$ of $V$ complementary to $W$. The dimension of $W'$ is $n-1$. Let's introduce to $V$ a inner product $(\cdot|\cdot)$ that fulfills the expected properties, in particular:
\begin{align*}
	(\bas_i|\bas_j) = \delta_{i,j} = \begin{cases}
		1 \text{, if $i=j$,}\\
		0 \text{, if $i \neq j$.}
	\end{cases}
\end{align*}
for the basis $(\bas_i)_{i=1}^n$. Then the basis $(\fbas_j)_{j=1}^{n-1}$ of $W'$ can be constructed by ensuring that each $\fbas_j$ is orthogonal to the basis vector of $W$, that is
\begin{align*}
	\left( \fbas_j \ \middle\vert \ \sum_{i=1}^n \bas_i \right) = 0
\end{align*}
for every $1 \leq j \leq n-1$. One can see that one such basis is found by for example choosing every basis vector $\fbas_j$ to be the difference of two sequent $\bas_i$, that is choosing the basis $(\bas_j - \bas_{j+1})_{j=1}^{n-1}$ for $W'$. This is called the \emph{Standard representation of $\Sym_n$.}

\begin{example}[Standard representation of $\Sym_3$]
	The matrices of the standard representation is found by studying the action inherited from $\Sym_3$ on the basis 
	\begin{align*}
		\begin{cases}
			\fbas_1 = \bas_1-\bas_2, \\
			\fbas_2 = \bas_2-\bas_3.
		\end{cases}
	\end{align*}
	For example, $(1,2)\cdot (\bas_1-\bas_2) = \bas_2-\bas_1 = -\fbas_1$ and $(1,2) \cdot (\bas_2-\bas_3) = \bas_1-\bas_3 = \fbas_1 + \fbas_2$, hence $\rho_{(1,2)} = \left(\begin{smallmatrix}
		-1 & 1 \\ 0 & 1
	\end{smallmatrix}\right)$. The matrices are presented in Table~\ref{table:standardS3}.
	
	\begin{table}[hbt!]
		\centering
		\begin{tabular}{r r r}
			$\rho_{(1)} = 
			\begin{pmatrix}
				1 & 0 \\
				0 & 1
			\end{pmatrix}$, & 
			$\rho_{(1,2,3)} = 
			\begin{pmatrix}
				0 & -1 &\\
				1 & -1
			\end{pmatrix}$, & 
			$\rho_{(1,3,2)} = 
			\begin{pmatrix}
				-1 & 1  \\
				-1 & 0
			\end{pmatrix}$, \\ & & \\
			$\rho_{(1,2)} = 
			\begin{pmatrix}
				-1 & 1 \\
				0  & 1
			\end{pmatrix}$, &
			$\rho_{(1,3)} = 
			\begin{pmatrix}
				-1 & 0 \\
				0 & -1 
			\end{pmatrix}$, &
			$\rho_{(2,3)} = 
			\begin{pmatrix}
				1 & 0 \\
				1 & -1 
			\end{pmatrix}$.
		\end{tabular}
		\caption{Matrices of the standard representations of $\Sym_3$}
		\label{table:standardS3}
	\end{table}
\end{example}

\subsection{Irreducible representations}

In Section~\ref{sect:subrepr} we found that a given representation can be divided into complementary subrepresentations, now we introduce the notion of an ``indivisible'' representation.

\begin{definition}[Irreducible representation]
	If there are no proper and non-trivial $G$-invariant subspaces of a representation $V$, then it is said to be \textit{irreducible}. 
\end{definition}

We have already met a few of these.

\begin{example}[Degree 1 representations are irreducible]\cite[Example 1.4.2.]{Sagan}
	A vector space of dimension 1 has no other subspace other than itself and the zero space, thus it is irreducible. Hence the trivial representation of degree 1 of any group and the alternating representation of $\Sym_n$ discussed in Section~\ref{sect:basicreprs} and the degree 1 representations of $\Cyc_n$ found in Section~\ref{sect:deg1cycn} are irreducible representations of their respective groups.
\end{example}

The results from Proposition~\ref{prop:complsubrepr} invites the notion of \textit{complete reducibility} of an arbitrary representation. This is presented in the following theorem.

\begin{theorem}[Maschke's theorem]\label{thm:maschkes}
	Let $G$ be a finite group and let $V$ be any representation space of $G$. Then $V$ is composed as a direct sum of a number of not necessarily distinct subrepresentations $W_i$ of $V$, or in other words,
	\begin{align*}
		V = \bigoplus_i W_i
	\end{align*}
	\begin{note}
		The direct sum can be expressed in terms of distinct subrepresentation $W_i$ of $V$ and their multiplicities $a_i$ as
		\begin{align*}
			V = \bigoplus_i a_i W_i
		\end{align*}
	\end{note}
\end{theorem}

\begin{example}[Decomposition of permutation representation of $\Sym_n$]
	Let $G = \Sym_n$, and let $P$ be the permutation representation space, $S$ be the standard representation space and $T$ be the trivial representation space, then $V$ is fully decomposed as 
	\begin{align*}
		P = T \oplus V.
	\end{align*}
\end{example}

%\begin{example}[Decomposition of regular representation of $\Sym_n$]
%	Again, let $G = \Sym_n$, and let $R$ be the regular representation space, 
%\end{example}

\begin{theorem}[Schur's lemma]
	content...
\end{theorem}

%	In Example~\ref{example:comptrivpermrepr}, a subrepresentation of the permutation representation complementary to the trivial representation was found. Show it is irreducible? 

%	Consider the action of a group $G$ on the basis $(\bas_i - \bas_{i+1})_{i=1}^n$. For a $g \in G$ we have
%	\begin{align*}
	%		g \cdot (\bas_i - \bas_{i+1}) &= \bas_{g(i)} - \bas_{g(i+1)}
	%	\end{align*}
%	
%	\begin{example}
	%		hello
	%	\end{example}
	


\clearpage{\thispagestyle{empty}}
\section{Character Theory}

This section is based on \cite[Ch.2.]{Serre}, as well as \cite[Sect.2.2.]{FultonHarris}.

Let $V$ be a vector space with basis $(\bas_i)_{i=1}^n$ and let $\rho: G \rightarrow \GL(V)$ be a representation. Now, for each $g \in G$, define a function $\chi_{\rho}: G \rightarrow \CC$ to be the trace of the matrix of $g$ in $V$. In other words $\chi_{\rho}(g) = \Tr(\rho_g)$. This function is called the \emph{character} of a representation.

We recall from linear algebra the concept of the trace of a matrix $(a_{ij})_{n \times n}$. It is the sum along the diagonal,
\begin{align*}
	\Tr(a_{ij})= \sum_{i=1}^{n} a_{ii} = a_{11} + \dots + a_{nn},
\end{align*}
and is also the sum of the eigen values of the matrix and is independent of the basis chosen. Why is it useful when studying representations?

\subsection{Basic properties of characters}

\begin{proposition}\cite[Prop.2.1.]{Serre}\label{prop:charidinvconj}
	The trace function \emph{characterizes} the representation in some useful ways:
	\begin{enumerate}
		\item[i)] The character of the identity element of $G$ is the degree of the representation, $\chi(e) = \dim V$.
		\item[ii)] The character of an element in $G$ is the complex conjugate of the character of the inverse element, $\chi(g^{-1}) = \overline{\chi(g)}$.
		\item[iii)] The character is constant under conjugation, if $g,h\in G$ are conjugate it implies $\chi(g) = \chi(h)$.
	\end{enumerate}
\end{proposition}
\begin{proof}
	\begin{enumerate}
		\item[i)] By Propostion~\ref{prop:homoidinv}, $\rho_e$ is the $n \times n$ identity matrix, where $n = \dim V$, hence $\chi(e) = \Tr \id = n = \dim V$.
		\item[ii)] Let $g \in G$ and let $g^{-1}$ be its inverse. The character of $g^{-1}$ is 
		\begin{align*}
			\Tr(\rho_{g^{-1}}) = \Tr(\rho_g^{-1})
		\end{align*}
		which is the sum of the eigen values of $\rho_g^{-1}$. We are free to choose an orthonormal basis for $\rho_g$, then $\rho_g$ is a normal matrix. Now, if $\lambda$ is an eigenvalue of $\rho_g$, then $1/\lambda$ is an eigen value of $\rho_g^{-1}$, \textit{TO BE COMPLETED}

		\item[iii)] It is well known that the trace is conserved under conjugation. \textit{CITATION, HOLST?} \qedhere
	\end{enumerate}
\end{proof}

Now we properly introduce the character.

\begin{definition}\label{def:char}
	Let $V$ be a representation space of a group $G$ of dimension $n$ and let $\rho_g = (r_{ij})$ be the matrix representation of a $g \in G$. Then the \emph{character} of $g$ in $V$ is defined to be the trace of the matrix representation, that is 
	\begin{align*}
		\chi_V(g) := \Tr \rho_g = \sum_{i=1}^{n} r_{ii}.
	\end{align*}
\end{definition}


\begin{notation}
	A ``character vector'', simply called the group character, of $G$ can be defined as the tuple containing the character of every element of $G$, ie. $\chi = (\chi(g))_{g \in G}$.
\end{notation}

\begin{example}
	The trivial character is $(1, \dots, 1)$ for any group.
\end{example}

\begin{example}
	The permutation representation was described in Example~\ref{ex:permS3}. The character of $\Sym_3$ in the permutation representation is $(3, 0, 0, 1, 1, 1)$.
\end{example}

\begin{remark}
	We note that since the trace is fixed under conjugation, any elements of a group that are conjugate have the same character, ie. for a conjugate pair $g,g' \in G$ there exists an $h \in G$ such that $g' = hgh^{-1}$, and their characters are
	\begin{align*}
		\chi(h) = \chi(gag^{-1}) = \chi(g).
	\end{align*}
\end{remark}

\begin{notation}
	If the group $G$ is partitioned into the conjugacy classes $[k_1], [k_2], \dots, [k_l]$, then the group character can be abbreviated to contain one representative $k_i$ from every class $[k_i]$, ie. $\chi= (\chi(k_1), \chi(k_2), \dots, \chi(k_l))$.
\end{notation}

\begin{example}
	The group $\Sym_3$ is partitioned into the classes $[(1)], [(1,2)]$ and $[(1,2,3)]$, and the group character of the permutation representation can be written as $(3,0,1)$.
\end{example}

Moving on to the toolbox presented in Section~\ref{sect:tensorrepr}, we propose the following:

\begin{proposition}\label{prop:charplustimes}
	Let $V$ and $W$ be representation spaces of a group $G$. For a $g \in G$, let $\chi_V(g)$ and $\chi_W(g)$ be its characters in those representations. Then we have,
	\begin{enumerate}
		\item[i)] The character in $V \oplus W$ is $\chi_V+\chi_W$.
		\item[ii)] The character in $V \otimes W$ is $\chi_V \cdot \chi_W$.
		\item[iii)] The characters of the symmetric and exterior squares are \begin{align*}
			\chi_{\SymP^2 V}(g) = \frac{1}{2}(\chi_V(g)^2 + \chi_V(g^2)) \quad \text{and} \quad \chi_{\AltP^2 V}(g) = \frac{1}{2}(\chi_V(g)^2 - \chi_V(g^2)),
		\end{align*} compatible with $V \otimes V = \SymP^2 V \oplus \AltP^2 V$.
	\end{enumerate}
\end{proposition}
\begin{proof}
	Statement i) and ii) are consequences of Section~\ref{sect:tensoralgebra}. 
	\begin{enumerate}
		\item[i)] The trace of $\rho_g^V \oplus \rho_g^W$ is clearly the sum of the traces of $\rho_g^V$ and $\rho_g^W$.
		\item[ii)] Likewise, the trace of $\rho_g^V \otimes \rho_g^W$ is found to be the product of the traces of $\rho_g^V$ and $\rho_g^W$.
		\item[iii)] \textit{To be written...} \qedhere
	\end{enumerate}
\end{proof}

\subsection{Orthogonality relations}

Continuing the discussion of Section~\ref{sect:irredreprs}, we let $V$ be an arbitrary representation and $W$ be an irreducible representation of some group $G$. 

We denote by $\Hom (V,W)$ the set of all homomorphisms from $V$ to $W$, and by the superscript $(\cdot) ^G$ we denote a subset which is fixed by $G$, for example 
\begin{align}\label{eq:defVG}
	V^G = \left\lbrace \vvec \in V \middle\vert g\cdot \vvec = \vvec, \ \forall g \in G \right\rbrace
\end{align} 
and clearly, $\Hom(V,W)^G$ is the set of all $G$-linear maps $V \rightarrow W$.

\begin{conjecture}
	The multiplicity of $W$ in $V$ is $\dim \Hom(W,V)^G$, that is the number of $G$-linear maps from $W$ to $V$, ie. \begin{align*}
		\dim \Hom(W,V)^G = |\{\text{$G$-linear maps $W \rightarrow V$}\}|.
	\end{align*}
\end{conjecture}

Conversely, if $W$ is arbitrary and $V$ is irreducible, then $\dim \Hom (V,W)^G$ is the multiplicity of $V$ in $W$. If both $V$ and $W$ are irreducible, then by Schur's Lemma we must have that
\begin{align}\label{eq:dimhomVWG}
	\dim \Hom (V,W)^G = \begin{cases}
		1, \text{ if $V$ and $W$ are isomorphic, and} \\
		0, \text{ else.}
	\end{cases}
\end{align}

Now, let $V^G$ be defined as the fixed points of $V$ under the action of $G$, as defined in Equation~\ref{eq:defVG}. Note that a map\footnote{Such a map is called an \textit{endomorphism of $V^G$} and we set $\End(V^G) := \Hom(V^G, V^G)$.} $\varphi: V^G \rightarrow V^G$ is not just a representation of $G$, but also a trivial representation by definition, then the number of all such maps, denoted $\dim \End(V^G)$ is the multiplicity of the trivial representation in $V$. 

The map associated with any $g$ is not generally a $G$-linear map, however as we have previously seen, we can construct such a map by averaging over $G$, that is we define a map $\varphi: V \rightarrow V$ by
\begin{align*}
	\varphi = \frac{1}{|G|} \sum_{g \in G} g,
\end{align*}
which is $G$-linear since for some $h \in G$, \begin{align*}
	h \cdot \varphi(h^{-1} \cdot \vvec) &= \frac{1}{|G|} \sum_{g \in G} (hgh^{-1}) \cdot \vvec \\
	&= \varphi(\vvec)
\end{align*} for any $\vvec \in V$, also it is projection of $V$ onto $V^G$ PROOF (image of $\varphi$ is $V^G$ and $\varphi:V^G \rightarrow V^G$). 

Now, the multiplicity of the trivial representation in $V$ is the trace of $\varphi$ WHY??, which is
\begin{align*}
	\dim V^G = \Tr \varphi &= \Tr \left( \frac{1}{|G|} \sum_{g \in G} g \right) \\
	&= \frac{1}{|G|} \sum_{g \in G} \Tr g & \text{(Trace is linear operator)} \\
	&= \frac{1}{|G|} \sum_{g \in G} \chi_{V}(g) & \text{(Def. of character)},
\end{align*}
hence \begin{align}\label{eq:multVG}
	\dim V^G = \frac{1}{|G|} \sum_g \chi_V(g).
\end{align}

Now, let's apply Equation~\ref~{eq:multVG} to the set of all maps $V \rightarrow W$, ie. $V$ becomes $\Hom (V,W)$, where both $V$ and $W$ are irreducible representations of $G$. Earlier (Section~\ref{sect:tensorrepr}), $\Hom(V,W)$ was identified with $V^* \otimes W$, hence the character of $\Hom(V,W)$ is $\overline{\chi_V(g)}\cdot\chi_W(g)$, then the multiplicity of the trivial representation in $\Hom(V,W)$ is by Equations~\ref{eq:dimhomVWG} and~\ref{eq:multVG}:
\begin{align*}
	\dim \Hom(V,W)^G = \frac{1}{|G|} \sum_g \overline{\chi_V(g)}\cdot\chi_W(g) = \begin{cases}
		1, \text{ if $V \cong W$, and} \\
		0, \text{ if $V \neq W$.}
	\end{cases}
\end{align*}

We have arrived at an expression that looks suspiciously familiar, letting the two characters be complex-valued vectors in some vector space, we have found a \textit{scalar product of characters}.

\begin{definition}[Scalar product of characters]
	Let $\varphi$ and $\psi$ be the group characters of a group $G$, then we define the scalar product of characters as
	\begin{align*}
		(\varphi | \psi) = \frac{1}{|G|} \sum_{g \in G} \overline{\varphi(g)} {\psi(g)}.
	\end{align*}
	This is an inner product since it is linear in $\varphi$, semilinear in $\psi$, as well is $(\varphi|\varphi) > 0 $ for any $\varphi \neq \0$ and $(\varphi|\varphi)=0$ if and only if $\varphi= \0$. PROVE THESE PROPERTIES.
\end{definition}

%\begin{notation}
%	The inner product of a character with itself is sometimes lazily denoted by \begin{align*}
%		|\varphi| := (\varphi|\varphi).
%	\end{align*}
%\end{notation}

We have also shown the following property on irreducible representations:

\begin{theorem}[Irreducibility criterion]\label{thm:irredcrit}
	The characters of irreducible representations are orthonormal. In other words, let $\chi \neq \psi$ be characters of irreducible representations of $G$ (called irreducible characters), then we have that 
	\begin{enumerate}
		\item[i)] $(\chi|\chi)= 1$, and
		\item[ii)] $(\chi|\psi) = 0$.
	\end{enumerate}
\end{theorem}

Now, let $V$ be the direct sum of irreducible representations $W_i$ of $G$ such that
\begin{align*}
	V = \bigoplus_i W_i.
\end{align*}
If $\chi_i$ is the character of $W_i$, then by Proposition~\ref{prop:charplustimes} the character of $V$ is
\begin{align*}
	\varphi = \sum_i \chi_i.
\end{align*}

Let $\chi$ be the character of an irreducible representation of $G$, then we have that
\begin{align*}
	(\varphi|\chi) = \sum_i (\chi_i|\chi).
\end{align*}
Since $\chi$ and all of the $\chi_i$ are irreducible characters, all of the inner products $(\chi_i|\chi)$ is either 1 or 0, depending on if $\chi$ and $\chi_i$ are of isomorphic representations. Hence, $(\varphi|\chi)$ will return the number of occurences of $\chi$ in $\varphi$. This number is called the \textit{multiplicity} of $\chi$ in $\varphi$.

\begin{remark}
	This also means that the composition of a representation $\varphi$ into a direct sum of irreducible subrepresentations $\chi_i$ is unique. \textit{Something on two compositions of the same $V$ must be identical.}
\end{remark}

\begin{remark}
	The converse is also true, if two representations have the same character, then they are isomorphic since they contain the same irreducible representations.
\end{remark}

\begin{note}
	Let $V$ be a representation with the composition 
	\begin{align*}
		V = \bigoplus_{i=1}^k m_i W_i = m_1W_1 \oplus m_2W_2 \oplus \dots,
	\end{align*}
	where the $m_i$ are the multiplicities of the irreducible $W_i$. If the character of $W_i$ is denoted by $\chi_i$, then the character of $V$ is
	\begin{align*}
		\varphi = \sum_{i=1}^k m_i \chi_i,
	\end{align*}
	where $m_i = (\varphi|\chi_i)$.	Taking the inner product of $\varphi$ with itself we have
	\begin{align}\label{eq:innerprodself}
		(\varphi|\varphi) = \sum_{i=1}^k m_i^2.
	\end{align}
\end{note}

\begin{theorem}
	The character $\varphi$ is irreducible if and only if $(\varphi|\varphi) = 1$.
\end{theorem}
\begin{proof}
	If $\varphi$ is irreducible, then by Theorem~\ref{thm:irredcrit} we have that $(\varphi|\varphi) = 1$. For the converse statement, if the sum in Equation~\ref{eq:innerprodself} is equal to one, we must have that only one of the $m_i$ is equal to 1, and the rest are equal to 0. Then $V$ is composed of only \textit{one} irreducible character, hence it is irreducible.
\end{proof}

%\textit{Something on the characters are basis of class function space, hence the number of irreducible representations of a group is at most the number of conjugacy classes in the group.}

Using these criteria and relations as a character calculus toolbox, we can study some groups, determine all of their irreducible representations and decompose any arbitrary representation.

\subsection{Decomposition of the regular representation}

Some representations where painstakingly found in the regular representation (the trivial in any group, along with the alternating in the symmetric group), however in general the following hold for the regular representation:

For a $g \in G$, it will act on a basis vector $\bas_h$ of the regular representation space $V$ by
\begin{align*}
	g \cdot \bas_{h} = \bas_{gh}
\end{align*}
and the resulting matrix $\rho_g$ can then be constructed by studying the action of this $g \in G$ on every $\bas_h$, where $h \in G$. However, the trace of this matrix, ie. the character $\varphi(g)$, only depends on the values on the diagonal, which are the fixed points under the action of $g$. What this means is that the $h$th column will have an 1 in the $h$ row if and only if $gh = h$, which holds only if $g = e$ since only $e$ leaves every other element of $G$ fixed, hence $\rho_g$ is a $|G| \times |G|$ permuation matrix with trace
\begin{align*}
	\varphi(g) = \begin{cases}
		|G| \text{, if $g = e$,} \\
		0 \text{, otherwise,}
	\end{cases}
\end{align*}
and the character of the regular representation is 
\begin{align*}
	\varphi = (|G|, 0, \dots, 0).
\end{align*}

Now taking the inner product of it with itself we have,
\begin{align*}
	(\varphi|\varphi) = \frac{1}{|G|} \sum_{g \in G} \varphi(g)\overline{\varphi(g)},
\end{align*}
all summands vanish except $g = e$, and since $\varphi(e) = |G|$ we have that
\begin{align*}
	(\varphi|\varphi) = \frac{1}{|G|} |G|^2 = |G|,
\end{align*}
hence by Equation~\ref{eq:innerprodself}, the square sum of the multiplicities of the regular representation is the order of the group. Also, assume that ${W_i}$ is the family of irreducible representations of a group with respective characters $\chi_i$, then the regular representation is a direct sum of these with (not-necessarily non-zero) multiplicities $m_i$, then we have that the multiplicity of some irreducible $\chi_j$ in $\varphi$ is
\begin{align*}
	(\varphi|\chi_j) = \sum_i (\chi_i|\chi_j) = \sum_i \frac{1}{|G|} \sum_{g \in G} \chi_i(g) \overline{\chi_j(g)}.
\end{align*}
The $\chi_i$ vanish for all $g \neq e$, 
\begin{align*}
	= \sum_i \frac{1}{|G|} \chi_i(e) \overline{\chi_j(e)} = \sum_i \frac{1}{|G|}|G|\chi_j(e) = \sum_i \chi_j(e)
\end{align*}\marginnote{some mathematical error here}
hence the multiplicity of an irreducible representation inside the regular representation is the degree of that irreducible.

We arrive at:
\begin{theorem}
	The regular representation $V$ of a group is decomposed in term of the irreducible representations of the group $W_i$ as
	\begin{align*}
		V = \bigoplus_i \dim (W_i) W_i.
	\end{align*}
\end{theorem}

\begin{proof}
	\textit{To be written...} \textit{HIGH PRIO}
\end{proof}

%Recall from linear algebra that two similar matrices have the same trace \cite[Thm 5.4.]{Holst} and that two matrices $X$ and $Y$ are similar if there exists an invertible matrix $S$ such that $S X S^{-1} = Y$. Also recall that the trace is the sum along the diagonal of the matrix, as well as the sum of the eigen values of the matrix.
%
%Something on change of basis - trace is basis independent.
%
%The trace seem related to the conjugation of elements of $\GL(V)$, motivating our use of it as a ``metric/measure'' to compare different representations. 
%
%%\begin{definition}\label{def:char}
%%	 Let $V$ be a representation space of a group $G$ of dimension $n$ and let $\rho_g = (r_{ij})$ be the matrix representation of a $g \in G$. Then the \emph{character} of $g$ in $V$ is defined to be the trace of the matrix representation, that is 
%%	 \begin{align*}
%%	 	\chi_V(g) := \Tr \rho_g = \sum_{i=1}^{n} r_{ii}.
%%	 \end{align*}
%%\end{definition}
%
%Whenever $V$ or $\rho$ are contextually obvious they may be omitted, ie. $\chi = \Tr g$. We have already inadvertently studied the characters of some representations:
%
%\begin{example}[Characters of degree 1 representations]
%	Evident from Definition~\ref{def:char}, the character of a degree 1 representation is the same as the representation, since a $1 \times 1$-matrix is usually identified with its only element, ie. $\chi_g = \rho_g$, hence the Tables~\ref{table:Cyc3},~\ref{tbl:cyc4} and \ref{table:Cyc5} are the character tables of $\Cyc_3$, $\Cyc_4$ and $\Cyc_5$ respectively.
%	
%	Further more, the character of the trivial representation of a group is 1 for every element of the group and the character of the alternating representation of the symmetric group is $+1$ for an even permutation and $-1$ for an odd one.
%\end{example}
%
%\begin{example}[Character of the identity element]
%	By Equation~\ref{eq:rhoIdAndInv}, the identity element of any group is mapped to the identity transformation of the representation $V$, which is the identity matrix, hence $\chi(e) = \dim V$.
%\end{example}
%
%\begin{example}[Character of the permutation representation of $\Sym_n$]
%	The matrix representation of a $\sigma \in \Sym_n$ has a 1 in the diagonal if and only if that row corresponds to a fixed point, hence the character of the permutation representation is the same as the number of fixed points of that conjugacy class.
%\end{example}
%Something on class functions
%
%Something on basis, matrix representation
%
%Projection formulae
%
%character tables 
%
%Sometging on the regular representation.
{\allowdisplaybreaks\subsection{Examples}

%\textit{REWRITE THIS SECTION AROUND THE STUDY OF THE REGULAR REPRESENTATION.}

\begin{notation}
	For this section, a representation space may be denoted by its designation in the subscript, that is a trivial representation space is denoted 
	by $V_\Triv$, a permutation by $V_\Perm$, etc.
\end{notation}

\subsubsection{Characters of $\Cyc_n$}

\begin{example}[Characters of $\Cyc_n$]
	By Equation~\ref{eq:regcharacter}, the character of the regular representation of $\Cyc_n$ is $\chi_\Reg = (n, 0, \dots)$. In section \ref{sect:deg1cycn} we described $n$ irreducible representations of degree 1 of  $\Cyc_n$, and by Theorem~\ref{thm:regreprmult}, we have then found every irreducible representation of $\Cyc_n$ since $\sum_{i=1}^n 1^2 = n$.
	
	Remember that the trace of a $1 \times 1$ matrix is its only element, therefore the Tables~\ref{table:Cyc3}, \ref{tbl:cyc4} and \ref{table:Cyc5} are the character tables of respectively $\Cyc_3$, $\Cyc_4$ and $\Cyc_5$.
\end{example}

\subsubsection{Characters of $\Sym_n$}
	
\begin{example}[Characters of $\Sym_3$]
	Returning to Table~\ref{table:chartableSym3}, we have found three irreducible representations of $\Sym_3$. Their degrees are 1, 1 and 2, and the square sum of the degrees are $1 + 1 + 4 = 6$ which is the order of $\Sym_3$, hence by Corollary~\ref{cor:sqsumirreds}, we have found every irreducible representation of $\Sym_3$. 
	
	Now we decompose the tensor powers of the standard representation.
	
	\paragraph{Decomposition of $V_\Stan \otimes V_\Stan$.} Its character is $(4,0,1)$, which by a quick glance on the character table is seen to be the sum of all irreducibles, hence \begin{align*}
		V_\Stan \otimes V_\Stan = V_\Triv \oplus V_\Alt \oplus V_\Stan.
	\end{align*}
	
%	\paragraph{Decomposition of $\SymP^2S$ and $\AltP^2S$.} The character of $\SymP^2S$ is 
%	\begin{align*}
%		\frac{1}{2}\left[ (4,0,1) - (2,2,-1) \right] = (1,-1,1) = \chi_\Alt,
%	\end{align*} hence the symmetric square of the standard representation is isomorphic to the alternating representation. 
%	
%	The character of $\AltP^2S$ is likewise 
%	\begin{align*}
%		\frac{1}{2}\left[ (4,0,1) - (2,2,-1) \right] = (3,1,0) = \chi_\Perm, 
%	\end{align*}
%	hence the exterior square of the standard representation is isomorphic to the permutation representation, compatible with $S \otimes S = \SymP^2S \oplus \AltP^2S$.
	
	\paragraph{Decomposition of $V_\Stan^{\otimes n}$.}\cite[Exercise 2.7.]{FultonHarris} To find the decomposition of larger tensor powers, we study the character $\chi_\Stan^n = (2^n, 0, (-1)^n)$ and take the inner product of it with the irreducibles and find:
	\begin{align*}
		({\chi_\Triv}|\chi_\Stan^n) = ({\chi_\Alt}|\chi_\Stan^n) &= \frac{1}{6}\left(2^n + (-1)^n\right), \quad \text{and}\\
		({\chi_\Stan}|\chi_\Stan^n) &= \frac{1}{6}\left(2^{n+1} + (-1)^{n+1}\right), 
	\end{align*}
	hence \begin{align*}
		V_\Stan^{\otimes n} = a_n (V_\Triv \oplus V_\Alt) \oplus a_{n+1} V_\Stan, \quad \text{where} \quad a_n = \frac{1}{6}\left(2^n + (-1)^n\right).
	\end{align*} 
\end{example}

	\paragraph{} We arrive at a ``complete'' character table for $\Sym_3$, presented in Table~\ref{table:completecharS3}.
	
	\begin{table}[hbt!]
		\centering
		
		\begin{tabular}{c | c c c}
			$S_3$         & $[(1)]$   & $[(1,2)]$  & $[(1,2,3)]$ \\
			$\Ss|[\sigma]|$    & $\Ss 1$ & $\Ss 3$ & $\Ss 2$ \\ \hline
			$\chi_\Triv$       & 1       & 1       & 1       \\
			$\chi_\Alt$       & 1       & -1      & 1       \\
			$\chi_\Stan$       & 2       & 0       & -1      \\ \hline\hline
			$\chi_\Perm$       & 3       & 1       & 0       \\
			$\chi_\Stan^2$ & 4       & 0       & 1       \\
			$\chi_\Reg$       & 6       & 0       & 0       		\end{tabular}
			
		\caption{Complete character table of $S_3$. The representations above the doublestruck line are irreducibles, and those below are composed.}
		\label{table:completecharS3}
	\end{table}
%\end{example}

\begin{example}[Characters of $\Sym_4$]
	This section follows the same methods from the previous section on $\Sym_3$. So far, we have that the trivial, the alternating and the standard representations are irreducibles, see Table~\ref{table:chartableSym4}. Their respective degrees are 1, 1 and 3, with the square sum $1+1+9 = 11$, which is less than $|\Sym_4|=24$, hence by Corollary~\ref{cor:sqsumirreds} we are expected to find additional irreducible representations of $\Sym_4$. We will find them while constructing and decomposing new representations.
	
	\paragraph{Decomposition of $V_\Alt \otimes V_\Stan$.} Consider an ``alternating version of the standard representation'', denoted with the space $V'_\Stan := V_\Alt \otimes V_\Stan$ and the character $\chi'_\Stan := \chi_\Alt\chi_\Stan = (3,-1,0,1,-1)$. It clearly has the same inner product with itself as $\chi_\Stan$, hence it is also irreducible. It increments the square sum of degrees to $11 + 3^2 = 20$.
	
	\paragraph{Decomposition of $V_\Stan \otimes V_\Stan$.} To find the next irreducible, we study the tensor square of the standard representation. Its character is $\chi_\Stan^2 = (9,1,0,1,1)$ and calculations will show that
	\begin{align*}
		(\chi_\Triv | \chi_\Stan^2) & = \frac{1}{24}(9 + 6 + 0 + 6 + 3) = 1, \\
		(\chi_\Alt | \chi_\Stan^2)  &= \frac{1}{24}(9-6+0-6+2) = 0, \\
		(\chi_\Stan | \chi_\Stan^2)  &= \frac{1}{24}(27+6+0-6-3) = 1, \quad \text{and} \\
		(\chi_\Alt\chi_\Stan | \chi_\Stan^2)  &= \frac{1}{24}(27-6+0+6-3) = 1.
	\end{align*}
	However these subrepresentations of non-zero multiplicity are of degrees 1, 3, and 3, and their direct sum is of degree 7, which is less than the dimension of $V_\Stan \otimes V_\Stan$, hence by Proposition~\ref{prop:complsubrepr} there is another (not necessarily irreducible) subrepresentation of degree 2 complementing them. Denoting this representation by $W$, its character is $\chi_W = \chi_\Stan^2 - \chi_\Triv - \chi_\Stan - \chi'_\Stan = (2,0,-1,0,2)$, which is found to be such that 
	\begin{align*}
		(\chi_W|\chi_W) = \frac{1}{24}(4 + 0 + 8 + 0 + 12) = 1,
	\end{align*}
	hence $W$ is irreducible and the tensor square of the standard representation is decomposed to 
	\begin{align*}
		V_\Stan \otimes V_\Stan = V_\Triv \oplus W \oplus V_\Stan \oplus V'_\Stan.
	\end{align*}
	We have now found all five irreducible representations of $\Sym_4$, since $1^2 + 1^2 +  2^2 + 3^2 + 3^2 = 24$, the order of $\Sym_4$.
%	\paragraph{Decomposition of $\SymP^2S$ and $\AltP^2S$.} The character of $\SymP^2S$ is 
%	\begin{align*}
%		\frac{1}{2}\left[ (9,1,0,1,1)-(3,3,0,-1,3) \right] = (3,-1,0,1,-1)
%	\end{align*}
%	which is the character of $S'$. Likewise, the character of $\AltP^2S$ is $(6,2,0,0,2)$ which is found to be the sum of the characters of $T, V$ and $S$, hence $\AltP^2S = T \oplus V \oplus S$.
	\paragraph{Decomposition of the $n$th tensor power of $W$.}	The character of $W^{\otimes n}$ is
	\begin{align*}
		\chi_{W}^n = (2^n, 0, (-1)^n, 0, 2^n).
	\end{align*}
	
	Its inner product with the irreducible characters are found to be
	\begin{align*}
		(\chi_\Triv| \chi_{W}^n) = (\chi_\Alt | \chi_{W}^n) &= \frac{1}{6} \left(2^n + 2 \cdot (-1)^n\right), \\
		(\chi_W | \chi_{W}^n) &= \frac{1}{6}\left(2^{n+1} + 2 \cdot (-1)^{n+1}\right), \text{ and}\\
		(\chi_\Stan | \chi_{W}^n) = (\chi'_\Stan | \chi_{W}^n) &= 0,
	\end{align*}
	hence the decomposition is
	\begin{align*}
		W^{\otimes n} = a_n (T \oplus A) \oplus a_{n+1} W, \quad \text{where} \quad a_n := \frac{1}{6}\left( 2^n + 2 \cdot (-1)^n \right).
	\end{align*}
	
	\paragraph{Decomposition of the $n$th tensor power of $V_\Stan$.} The character of $V_\Stan^{\otimes n}$ is 
	\begin{align*}
		\chi_\Stan^n = (3^n, 1, 0, (-1)^n, (-1)^n),
	\end{align*}
	and after projecting it on the irreducible characters we find that
	\begin{align}\label{eq:nthpowerS}
		V_\Stan^{\otimes n} = a_n V_\Triv \oplus b_n V_\Alt \oplus c_n W \oplus a_{n+1} V_\Stan \oplus b_{n+1} V'_\Stan,
	\end{align}
	where 
\begin{align*}
		\begin{cases}
		a_n = \frac{1}{24}(3^n + 9\cdot(-1)^n+6), \\
		b_n = \frac{1}{24}(3^n - 3\cdot(-1)^n -6), \quad \text{and} \\
		c_n = \frac{1}{12}(3^n + 3\cdot(-1)^n).
	\end{cases}
\end{align*}
	
	\paragraph{Decomposition of the $n$th tensor powers of $V'_\Stan$.} Since $V'_\Stan = V_\Alt \otimes V_\Stan$, the character of $(V'_\Stan)^{\otimes n}$ is the $n$th power of the character $\chi_\Alt \chi_\Stan$, which is $\chi_\Alt^n\chi_\Stan^n$, hence $(V'_\Stan)^{\otimes n} = V_\Stan^{\otimes n} \otimes V_\Alt^{\otimes n}$.
	
	By Example~\ref{example:nthpoweralternating}, the $n$th tensor power of $V_\Alt$ is $V_\Triv$ if $n$ is even and $V_\Alt$ if $n$ is odd, so for even $n$, $(V'_\Stan)^{\otimes n}$ is isomorphic to $V_\Stan^{\otimes n}$, and for odd $n$, it is isomorphic to $ V_\Stan^{\otimes n} \otimes V_\Alt$. Tensor multiplying Equation~\ref{eq:nthpowerS} with $V_\Alt$, we have:
	\begin{align*}
		V_\Stan^{\otimes n} \otimes V_\Alt &= \Big(a_n V_\Triv \oplus b_n V_\Alt \oplus c_n W \oplus a_{n+1} V_\Stan \oplus b_{n+1} V'_\Stan\Big)\otimes V_\Alt \\
		&= b_n V_\Triv \oplus a_n V_\Alt \oplus c_n W \oplus b_{n+1} V_\Stan \oplus a_{n+1} V'_\Stan,
	\end{align*}
	since 
	\begin{align*}
		V_\Triv \otimes V_\Alt &= V_\Alt, \\
		V_\Alt \otimes V_\Alt &= V_\Triv, \\
		V_\Stan \otimes V_\Alt &= V'_\Stan, \quad \text{and} \\
		V'_\Stan \otimes V_\Alt &= V_\Stan \otimes V_\Alt^{\otimes 2} = V_\Stan.
	\end{align*}
	Then we have found the decomposition of all tensor powers of $V'_\Stan$: For even $n$, it is the same as $V_\Stan^{\otimes n}$, and lastly for odd $n$, $V_\Triv$ switch multiplicities with $V_\Alt$, and $V_\Stan$ switch with $V'_\Stan$.
	
	\paragraph{} The findings of the last few paragraphs (except those on larger tensor powers) are presented in table \ref{table:completecharS4}.
%	 \vspace{2cm}
%	\vfill
	
	\begin{table}[hbt!]
		\centering
		\begin{tabular}{c | c c c c c}
			$\Sym_4$         & $[(1)]$   & $[(1,2)]$  & $[(1,2,3)]$ & $[(1,2,3,4)]$ & $[(1,2)(3,4)]$ \\
			$\Ss|[\sigma]|$    & $\Ss 1$ & $\Ss 6$ & $\Ss 8$ & $\Ss 6$  & $\Ss 3$      \\ \hline
			$\chi_\Triv$       & 1       & 1       & 1       & 1        & 1            \\
			$\chi_\Alt$       & 1       & -1      & 1       & -1       & 1            \\
			$\chi_W$       & 2       & 0       & -1      & 0        & 2            \\
			$\chi_\Stan$       & 3       & 1       & 0       & -1       & -1           \\
			$\chi'_\Stan$ & 3       & -1      & 0       & 1        & -1           \\ \hline\hline
			$\chi_\Perm$       & 4       & 2       & 1       & 0        & 0            \\
			%$\chi_{\bigwedge^2S}$    & 6       & 2       & 0       & 0        & 2            & $\bigwedge^2S \cong T \oplus S \oplus V$                     \\
			$\chi_\Stan^2$ & 9       & 1       & 0       & 1        & 1            \\
%			$\chi_R$       & 24      & 0       & 0       & 0        & 0            & $R \cong T \oplus A \oplus 2V \oplus 3S \oplus 3S'$
		\end{tabular}
		\caption{Complete character table of $\Sym_4$. The representations above the doublestruck line are irreducibles, and those below are composed.}
		\label{table:completecharS4}
	\end{table}
\end{example}

%\textit{LOW PRIO: ADD CHARACTER THEORY ON KLEIN'S 4-GROUP AND DIHEDRAL GROUPS}

%\subsection{Characters of the Klein 4-group $V$}
%
%The Klein 4-group, denoted by\footnote{$V$ is for vierer, four in german.} $V$ is, along with $Z_4$, the only group of order 4. It can be described as the set
%\begin{align}
%	V = \{ e, x, y, xy \},
%\end{align}
%where every element is its own inverse and the product of any two distinct non-identity elements is the third one.
%It is an abelian group, which means that every element has its own conjugacy class (SOURCE).
%
%Since every group has a trivial representation $T$ in any vector space, so do $V$. The center of $V$ is trivial, so the regular representation $R$ has the character $(4,0,0,0)$ and is clearly not irreducible.
%
%If $V$ instead is considered as a subgroup of $S_4$, that is
%\begin{align}
%	V = \{ (1), (12)(34), (13)(24), (14)(23) \} < S_4,
%\end{align}
%then each entry in the character of the permutation representation $P$ is the number of fixed points of the associated element. Then one sees that the character of the permutation representation is also $(4,0,0,0)$, it isomorphic to the regular representation.
%
%From theorem \ref{thm:numberirrep} we expect three additional irreducible representations. Since the square sum over all representations of the character of $e$ must be equal to the order of the group (SOURCE), they must be of degree 1. We know that they must be parallel to themselves and orthogonal to $T$ and have a 1 as the character of $e$, so the three additional are alterations of $T$ with two of the ones exchanged with negative ones. Then $R$ is the direct sum of these four irreducibles. These findings are presented in table \ref{table:charV} as the representations $A, B$ and $C$.
%
%
%\begin{table}[hbt!]
%	\centering
%	\begin{tabular}{c | c c c c l}
%		$V$ & $e$ & $x$ & $y$ & $xy$ &             \\ \hline
%		$T$ & 1   & 1   & 1   & 1   &             \\
%		$A$ & 1   & -1  & -1  & 1   &             \\
%		$B$ & 1   & -1  & 1   & -1  &             \\
%		$C$ & 1   & 1   & -1  & -1  &             \\ \hline\hline
%		$R$ & 4   & 0   & 0   & 0   & %$R \cong P \cong T \oplus A \oplus B \oplus C$
%	\end{tabular}
%	\caption{Character table for $V$.}
%	\label{table:charV}
%\end{table}
%
%
%\subsection{Characters of the dihedral group of order 8, $\Dih_8$}
%
%$\Dih_8$ is a 2-Sylow subgroup of $S_4$. Restrict the character table of $S_4$ to $\Dih_8$. 
%
%$\Dih_8$ has its non-trivial                                                                                                                                                                                                                                                              center $Z(\Dih_8) = \{e, r^2\}$ as a normal subgroup. The quotient $\Dih_8/\{e, r^2\}$ is isomorphic to the Klein 4-group. The character table of the Klein 4-group is a ``subset'' of the character table of $\Dih_8$.
}

%\clearpage{\thispagestyle{empty}}
\section{Conclusion}

A conclusion will be written here.  \lipsum[8]

Further topics. \lipsum[9]

Alternate paths (modules) \lipsum[10]

%The section with references.

\clearpage{\thispagestyle{empty}}
%\bibliographystyle{alphaurl} %This style prints the urls of online sources.
\bibliographystyle{alpha}
\bibliography{bibFile}

% All references are to be stored in the file bibliographyFile.bib! 
% It is in the same folder as your main .tex file.
% It is using the bibtex syntax, and you can get these entries from google scholar, etc.
% Your supervisor should know,

%\bibliography{bibliographyFile} 

\end{document}
